\chapter{共同犯罪}

% Source PDF physical page 133; see audit/source_coverage.json.
\begin{EditionOpeningPanel}
共同犯罪在刑法总论中有两个最：地位最重要，难度最大。希望勇敢的你耐心学习，多看两遍。

法考主要考查两大板块。第一，四大共犯人角色：共同正犯、间接正犯、教唆犯、帮助犯。第二，四大角色与其他知识点的结合，主要考查共犯人之间“步调不一致”。
\end{EditionOpeningPanel}

\section{共同犯罪的分类与原理}

\subsection*{一、共同犯罪的分类}

\subsubsection*{（一）分工分类}

从分工上看，共同犯罪分为教唆犯、帮助犯、实行犯。例如，甲教唆乙丙共同盗窃丁家。乙提供钥匙，丙入户盗窃。甲是教唆犯，乙是帮助犯，丙是实行犯。教唆犯、帮助犯统称为狭义共犯（广义的共犯就是指共同犯罪人，包括了所有参与犯罪的人）。实行犯又简称正犯（正犯除了包括实行犯，还包括其他没有实行行为但起到重要作用的正犯，对此后文会讲解）。

实行犯与狭义共犯的区分标准是：\CJKunderwave{对法益的侵害方式不同}。实行犯对法益的侵害具有\CJKunderwave{直接性、支配性}。教唆犯、帮助犯对法益的侵害具有\CJKunderwave{间接性}，需要凭借实行犯才能发挥作用；教唆犯、帮助犯对法益的侵害缺乏支配性。

[提示】分析共同犯罪的案件，\CJKunderwave{正确的分析顺序是，先分析正犯（主要是实行犯），}\CJKunderwave{后分析共犯（教唆犯和帮助犯）。要树立以正犯为中心的观念}。

\InsertDiagram{figures/manual_p133_1.tex}

1.从数量上看，正犯包括单独正犯和共同正犯。例1，甲教唆乙抢劫银行。乙抢劫了银行。甲是教唆犯，乙是单独正犯。例2，甲、乙共谋一起抢劫银行，并且一起使用暴力抢劫了银行。甲、乙属于共同正犯（共同实行犯）。

2.从方式上看，正犯包括直接正犯和间接正犯。直接正犯是指亲自实施实行行为。间接正犯是指不亲自实施，利用他人实施，对他人有支配力。间接正犯对法益的侵害虽然

% Source PDF physical page 134; see audit/source_coverage.json.
不具有直接性，但具有支配性。例1，甲指使7岁小孩乙入户盗窃，乙照办。乙是甲的犯罪工具，甲是间接正犯。如果甲亲自盗窃，则甲是直接正犯。例2（2023年试题），下列说法是否正确：“间接正犯的对应概念是单独正犯”？答案：说法错误。

[提示]考试主要考查四大角色：共同正犯，间接正犯，教唆犯，帮助犯。也即，两个正犯，两个共犯。

3.共同犯罪中，如果只有共同正犯，缺少共犯（教唆犯、帮助犯），则被称为简单共同犯罪。如果存在共犯，则被称为复杂共同犯罪。

\InsertDiagram{figures/manual_p134_1.tex}

\subsubsection*{（二）作用分类}

从作用来看，共同犯罪分为主犯、从犯、胁从犯。主犯起主要作用。从犯起次要作用。胁从犯，是被胁迫参加犯罪。

\subsubsection*{（三）结构分类}

从结构上看，共同犯罪分为任意共同犯罪和必要共同犯罪。

1.任意的共同犯罪（任意共犯），是指一人能够单独实施的犯罪由二人以上共同实施的情形。例如，故意杀人罪、盗窃罪，既可以由一人实施，也可以由二人以上共同实施。

2.必要的共同犯罪（必要共犯），是指必须由二人以上共同实施的犯罪。例如，聚众斗殴罪，不可能由一人实施。必要共犯包括聚众共同犯罪、集团共同犯罪（在本讲最后一节有讲），也包括一部分“对向犯”。

[注意］对向犯，是指必须由二人实施，且二者的行为是相对或相向的。对向犯包括四种情形：

（1）双方的罪名、法定刑均相同。例如，重婚罪，代替考试罪（替考者与被替考者）。

（2）双方的罪名不同，法定刑相同。例如，出售、购买假币罪。出售者构成出售假币罪，购买者构成购买假币罪，但两罪的法定刑相同。

（3）双方的罪名、法定刑均不同。例如，行贿罪与受贿罪。老板甲向官员乙行贿10万元，乙接收。甲构成行贿罪，乙构成受贿罪。两罪的法定刑不同，后者重于前者。\footnote{注意：虽然在事实结构上，甲乙构成共同犯罪，甲构成乙的受贿罪的共犯（教唆犯、帮助犯），但是由于这种共犯行为已经被立法者规定为独立罪名（行贿罪），且配置了较轻的法定刑，因此，对甲只能以行贿罪加以认定，不能再认定为受贿罪的共犯。注意，这是就两者间的犯罪结构而言的。如果在甲乙之间有个中间人丙，穿针引线，促成权钱交易，那么丙既构成甲的行贿罪的共犯，也构成乙的受贿罪的共犯，还构成介绍贿赂罪，想象竞合，择一重罪论处。（2026年贪污贿赂犯罪司法解释二）}

% Source PDF physical page 135; see audit/source_coverage.json.
（4）只处罚一方的行为。例如，贩卖淫秽物品牟利罪，立法者只处罚贩卖者，不处罚购买者。\footnote{注意：由于立法者不处罚购买者，因此不能认定购买者构成贩卖淫秽物品牟利罪的共犯。但是，如果乙原本不是淫秽物品的贩卖者，甲教唆乙将私藏的淫秽物品卖给甲，乙由此答应照办，那么甲构成乙的贩卖淫秽物品牟利罪的教唆犯。}又如，贩卖毒品罪，立法者只处罚贩卖者，不处罚买来吸食的购买者。这种情形被称为片面（单方面）的对向犯。\footnote{提示：对向犯是否构成共同犯罪，需要具体情形具体分析，不能认为“对向犯一律构成共同犯罪”。例1，甲男有妻子，乙女单身，甲男谎称自己单身，与乙女结婚。甲男构成重婚罪，乙女无罪。甲乙不构成共同犯罪。例2，甲向官员乙行贿，将银行卡放到乙家沙发缝，乙不知情。甲构成行贿罪，乙不构成受贿罪。甲乙不构成共同犯罪。概言之，对向犯中，有些构成共同犯罪，属于必要共犯，有些不构成共同犯罪，不属于必要共犯。对此需要具体分析。必要共犯与对向犯是交叉重合关系，不是大集合与子集的关系。}

\subsection*{二、共同犯罪的原理}

\subsubsection*{（一）共同犯罪中“共同”的含义：相同性？一起干！}

“共同”一词有两个含义。一是形容词“相同的”，修饰名词。例如，“我们有共同的爱好”，是指有相同的爱好。二是副词“一起去”，修饰动词。例如，“我们共同投资做生意”，是指一起投资，一起做生意。

传统理论将共同犯罪的“共同”理解成形容词，将“犯罪”理解成名词，\CJKunderwave{将“共同}\CJKunderwave{犯罪”理解成“共同的犯罪”“相同的犯罪”}，认为二人是否构成共同犯罪，主要看二人的犯罪特征有无相同性，要求客观行为相同、主观故意相同、触犯罪名相同。然而，这是从形式特征来判断，是不妥当的。

正确理解应是，共同犯罪的“共同”是副词，“犯罪”是动词，\CJKunderwave{“共同犯罪”是指}\CJKunderwave{“共同去犯罪”，也即一起制造违法事实（法益侵害事实）}。概言之，\CJKunderwave{判断二人是否成立共}\CJKunderwave{同犯罪，应看违法事实的一起性，而非犯罪特征的相同性}。“一起性”的判断标准是二人是否具有意思联络，有意思联络表明是一起干。

例如，甲乙互不知情，同时向丙开枪，均打中丙，丙死亡。虽然甲乙的客观行为相同、主观故意相同、触犯罪名相同，但是，甲乙不是共同犯罪，因为甲乙没有意思联络，不属于共同制造违法事实，在所制造的违法事实上没有一起性（连带性），因此不构成共同犯罪。甲乙的这种犯罪现象被称为\CJKunderwave{同时犯}。

\subsubsection*{（二）共同犯罪中“犯罪”的含义：制造违法事实（干坏事）}

根据第二讲“犯罪构成”的讲解，“犯罪”概念要阶段化理解，第一个阶段的“犯罪”概念是“干坏事”（制造违法事实）。第二个阶段的“犯罪”概念是“干坏事并且是坏人”。

共同犯罪中的“犯罪”，是指第一个阶段的“犯罪”概念，也即“干坏事”，不要求是第二个阶段的“犯罪”概念，也即“干坏事并且是坏人”。

“共同犯罪”就是指，一起制造违法事实，也即“一起干坏事”。

第一，是不是“一起干”，看二人是否具有意思联络。

第二，“干坏事”指制造违法事实。

第三，成立共同犯罪，不要求行为人是“坏人（具有可谴责性的人）”，也即不要求

% Source PDF physical page 136; see audit/source_coverage.json.
行为人具有责任年龄、责任能力。共同犯罪是对事不对人。

例如，10岁的狗蛋为情所困，要杀小芳，指使老大爷为其望风。后狗蛋得逞。（老大爷望风案，2012年第9题的模型，2022年试题模型）。二人一起制造了违法事实。这种共同违法行为就是共同犯罪行为。到了排除责任（谴责）阶段，狗蛋不具有责任年龄，排除可谴责性，最终不负刑事责任。但是，老大爷最终应负刑事责任，构成故意杀人罪的帮助犯。若问故意杀人罪的实行犯在哪里，就是狗蛋，是“制造违法事实”意义上的“实行犯”。以上分析既是“九大本”法考教材的观点（2026年版刑法卷第59页），也是官方答案的观点。

[提示］苏联版四要件体系认为，共同犯罪的成立条件是，犯罪主体必须是具有责任年龄、责任能力的人，也即，成立共同犯罪的犯罪人必须是“坏人”（具有可谴责性的人）。由于故意杀人罪的责任年龄是12岁，因此，狗蛋不是合格的犯罪主体，直接无罪。如此，对老大爷只能单独处理。单独处理，就对老大爷无法定罪，出现处罚漏洞。对老大爷不能定帮助犯，是因为帮助犯不能单独存在。对老大爷也不能定间接正犯，是因为间接正犯要求对被利用人具有支配力，将其作为犯罪工具加以支配。老大爷没有指使支配狗蛋，老大爷是狗蛋的“马仔”。假如狗蛋是个好孩子，老大爷指使狗蛋杀害小芳，则老大爷是间接正犯。可见，苏联四要件体系会导致处罚漏洞。这是其不足之处。

[常考模型］违法事实阶段，都打勾；排除责任（谴责）阶段，一个打勾，一个打叉。

\InsertDiagram{figures/manual_table_p136_1.tex}

[总结】按照中国化改良版四要件体系，\CJKunderwave{共同犯罪的成立条件是：一起制造违法事实。}\CJKunderwave{分析共同犯罪的两句标语：“违法是一起的，责任是独立的}”。在“制造违法事实”阶段，看二人是否一起制造了违法事实（一起干坏事）。在排除责任（谴责）阶段，谴责谁、不能谴责谁（也即谁是坏人、谁不是坏人），独立进行。

[注意］罪名问题。罪名只是对犯罪事实的一种称谓。两个人构成共同犯罪，不要求最终定相同罪名，也不要求有个罪名来表达这个“共同犯罪”。这是因为，共同犯罪是指一起制造违法事实。对这个事实，有名称（罪名）来表达更好，没有名称也无所谓；不能因为没有名称，就否定共同犯罪的事实，不能因“名”而废“实”。

例如，甲给乙发微信：“一起找个女的，搞一下。”乙回复：“好”。二人抓住丙女，准备带到街角，此时被警察抓获。单独讯问甲，甲答：“为了抢劫”。单独讯问乙，乙答：“甲想强奸，我跟他一起干。”甲乙构成共同犯罪，因为一起制造违法事实（抓住丙女）。没有某个罪名能够表达这个共同犯罪，但无所谓。甲构成抢劫罪未遂，乙构成强奸罪未遂。

\section{共同正犯}

[问题】猎人甲乙以为前方是野猪，商量道：“我们一起开枪，打死野猪。”不慎打死了行人丙。能够查明，致命的只有一颗子弹，但无法查明致命子弹是谁打的。对甲乙该如何处理？（打野猪案）

% Source PDF physical page 137; see audit/source_coverage.json.
\CJKunderwave{共同正犯的成立条件}是：通过实行行为，一起制造违法事实。例如（抢银行案），甲乙共谋一起抢劫银行，甲持枪压制银行职员，乙用麻袋装钱。甲乙通过实行行为，一起制造违法事实，构成共同正犯。

\InsertDiagram{figures/manual_table_p137_1.tex}

\subsection*{一、客观要件：一起制造违法事实}

\subsubsection*{（一）“一起实行”与意思联络}

如何判断二人是否“一起”实施实行行为？看二人是否有意思联络。

1.双方均没有意思联络，成立\CJKunderwave{同时犯}，不成立共同犯罪。例如，甲乙互不知情，同时向丙开枪；甲击中丙，乙未击中。甲乙不成立共同犯罪，仅属于同时犯。甲构成故意杀人罪既遂，乙构成故意杀人罪未遂。

2.双方均有意思联络，成立共同犯罪。例如，在上述“抢银行案”中，甲乙共谋，相互有意思联络，表明二人是一起干，构成共同正犯。

\subsubsection*{（二）法律后果：“部分实行、全部负责”}

甲乙有意思联络，在制造违法事实上具有一起性，那么在“谁需要为违法结果负责”的问题上就产生了连带责任，也即“\CJKunderwave{部分实行、全部负责}”原则。这是指，即使只有一个实行犯制造了违法结果，全部正犯都要为此负责。其理由根据是，二者有意思联络，在心理上相互加持，形成合力，有相互的心理性贡献。

例如（2008年第52题），甲与乙基于意思联络，同时向王某开枪，甲击中，导致王某死亡。乙未击中。甲乙构成共同正犯，乙应对甲制造的既遂结果负责，因此，乙也构成故意杀人罪既遂。\footnote{“部分实行、全部负责”原则适用于共同正犯之间。教唆犯、帮助犯对实行犯制造的结果是否负责，关键看教唆犯、帮助犯是否引起、促进实行犯的违法结果。对此在后文的“教唆犯”“帮助犯”中讲解。}

[注意］“部分实行、全部负责”原则的理由是“甲乙构成共同正犯”。然而，这是形式理由。实质理由是，甲乙有意思联络，相互有心理性贡献，在制造违法结果上具有连带责任。这表明，“甲乙构成共同正犯”是个形式称谓。即使去掉这个形式称谓，也不影响“乙对甲制造的违法结果负责”这个结论。

[引申总结】犯罪“既遂”的含义有两种：一是，自己亲力亲为制造了既遂结果，这是自然事实意义上的既遂结果；二是，自已为共同犯罪人的既遂结果负责，这是法律评价意义上的既遂结果。

例1，甲乙一起越狱，相互帮忙，乙逃出去了，甲未能逃出，被抓。乙构成脱逃罪既遂，甲也构成既遂。有同学不理解：“甲没逃出去，为何也构成既遂？”这里有两个评价角度。从自然事实角度看，甲没逃出去，单论甲，的确是未遂。但是，从法律评价角度看，甲要对乙的既遂结果负责，所以甲构成既遂。这里“甲构成既遂”的意思，不是甲自己既遂，而是指甲要对乙的既遂负责。

例2，甲乙一起强奸丙女，乙奸淫成功，甲未奸淫成功，被抓。乙构成强奸罪既遂，甲也构成既遂。这里“甲构成既遂”的意思，不是甲自己既遂，而是指甲要对乙的既遂负责。

% Source PDF physical page 138; see audit/source_coverage.json.
[常考情形]无法查明的情形。

例1，甲乙基于意思联络，共同向丙开枪，欲杀丙，丙死亡，只有一枪致命，但无法查明致命的子弹是谁打的。由于甲乙构成共同正犯，根据“部分实行、全部负责”的原则，此时无须查明是谁打的，因为每个人均需对另一人制造的结果负责。因此，甲乙均构成故意杀人罪既遂。

例2，甲乙互不知情，在不同地点同时向丙开枪，欲杀丙，丙死亡，只有一枪致命，但无法查明致命的子弹是谁打的。由于甲乙不构成共同正犯，而属于同时犯，违法行为不具有一起性，不能适用“部分实行、全部负责”的原则。对此只能适用存疑时有利于被告原则，有利于被告的结论就是甲乙各自构成故意杀人罪未遂。

\subsection*{二、主观要件}

例1（同时犯），猎人甲乙没有意思联络，互不知情，在不同地点同时向“野猪”开枪，实际打死了丙。只有一颗子弹打中，无法查明致命的这颗子弹是谁打的。由于甲乙没有意思联络，不是共同犯罪，各自单独处理。对此只能适用存疑时有利于被告原则，甲乙均对死亡结果不用负责，均不构成过失致人死亡罪，均无罪。

例2，本节开头的“打野猪案”，猎人甲乙以为前方是野猪，说：“我们一起开枪，才能打死野猪！”不慎打死了丙。只有一颗子弹打中，无法查明致命的这颗子弹是谁打的。注意，该案与例1的同时犯不同。该案中，甲乙有意思联络，属于“一起干”，不是各自单干。对该案，有三种观点。请细心的同学耐心体会。

\subsubsection*{（一）第一种观点：要求故意且无罪说}

1.结论：甲乙都无罪。理由：刑法第25条规定：“共同犯罪是指二人以上共同敌意犯罪。二人以上共同过失犯罪，不以共同犯罪论处………”成立共同犯罪，要求“故意+故意”。甲和乙是共同过失犯罪，不构成共同犯罪，对甲乙只能单独处理。单独处理时，遇到无法查明的情形，适用存疑时有利于被告原则，结论是，甲乙对死亡结果都不用负责，因此都不构成过失致人死亡罪，只能作无罪处理。

2.不足：甲乙有意思联络，一起打死了人，却都无罪，不符合实质正义。因此，该说是少数说。\footnote{“要求故意且无罪说”有时也被称为“部分犯罪共同说”。}

\subsubsection*{（二）第二种观点：不要求故意且有罪说}

1.结论：甲乙构成过失致人死亡罪的共同犯罪。理由：二人要成立共同正犯，只要求具有意思联络、一起制造违法事实。至于在主观要件上，可以是故意，也可以是过失。过失犯罪可以构成共同正犯。基于此，甲乙有意思联络，一起制造了死亡结果，构成共同正犯。根据“部分实行，全部负责”原则，对于无法查明的情形，无需查明，二人对死亡结果都要负责。到了主观要件，二人对死亡结果持过失心理，那么二人均构成过失致人死亡罪。

2.不足：该说认为，“共同过失犯罪能够构成共同犯罪”；这种看法违反了我国刑法规定，不符合形式正义。因此，该说是少数说。\footnote{“不要求故意且有罪说”有时也被称为“行为共同说”。}

\subsubsection*{（三）第三种观点：要求故意且有罪说（结果归责说）}

1.结论：甲乙不成立共同犯罪，各自单独构成过失致人死亡罪。理由：

% Source PDF physical page 139; see audit/source_coverage.json.
如前文所述，“部分实行、全部负责”原则的理由是“甲乙构成共同正犯”。然而，这是形式理由。实质理由是，甲乙有意思联络，相互有心理性贡献，在制造违法结果上具有连带责任。这表明，“甲乙构成共同正犯”是个形式称谓，是个“名”。即使去掉这个形式称谓，也不影响“一方对另一方制造的违法结果负责”这个事实和结论。在此应注意“名”与“实”的关系。不能因“名”而废“实”，可以去“名”而保“实”。也即，可以不要“甲乙构成共同正犯”这个“名”，但要保留这个“实”，也即“甲乙有意思联络，相互有心理性贡献，在制造违法结果上具有连带责任，也即部分实行、全部负责”。换言之，即使认为“甲乙不构成共同正犯”，但只要甲乙有意思联络，就能启用“部分实行、全部负责”原则。

\InsertDiagram{figures/manual_p139_1.tex}

在“打野猪案”中，甲乙有意思联络，相互有心理性贡献，形成心理上的合力，一起制造违法事实，那么在制造违法结果上应负连带责任，也即可以启用“部分实行、全部负责”原则。基于此，无论谁打死了丙，另一方都为此负责。因此，甲乙均构成过失致人死亡罪。至于要不要称谓“甲乙是共同正犯”，这个“名”不重要。对此，可以按照法条规定，认为甲乙不构成共同正犯。

\CJKunderwave{要求故意且有罪说认为，甲乙不构成共同犯罪，但各自单独构成过失致人死亡罪}。该观点既符合实质正义，也符合形式正义，是法考界的多数说。该说也被称为结果归责说。\footnote{参见王钢：《刑法犯罪论核心知识与阶层式案例分析》，北京大学出版社2025年版，第154页。}

2.法考真题的证明轨迹

（1）2008年四川卷真题考查了“打野猪案”，官方答案采取了要求故意且无罪说，认为甲乙都无罪。

（2）到了2016年第7题，官方答案改变了，采取了要求故意且有罪说，认为甲乙都有罪。对于官方答案的选择，应采取从新原则。

例如（2016年第7题），甲、乙、丙共同故意伤害丁，丁死亡。经查明，甲、乙都使用铁棒，丙未使用任何凶器；户体上除一处致命伤外，再无其他伤害；可以肯定致命伤不是丙造成的，但不能确定是甲造成还是乙造成的。关于本案，下列哪一选项是正确的？

A.因致命伤不是丙造成的，户体上也没有其他伤害，故丙不成立故意伤害罪

B.对甲与乙虽能认定为故意伤害罪，但不能认定为故意伤害（致死）罪

C.甲、乙成立故意伤害（致死）罪，丙成立故意伤害罪但不属于伤害致死

D.认定甲、乙、丙均成立故意伤害（致死）罪，与存疑时有利于被告的原则并不矛盾

官方答案：D。故意伤害罪（致人死亡）这一结果加重犯中，致人死亡是指过失致人

% Source PDF physical page 140; see audit/source_coverage.json.
死亡。题中，甲乙丙构成故意伤害罪的共同正犯，对死亡结果持过失心理。但是，三人是否构成故意伤害罪（过失致人死亡）这一结果加重犯？焦点问题是，甲乙丙对“过失致人死亡”是否负责？

“要求故意且无罪说”认为，针对“过失致人死亡”，甲乙不构成共同正犯，各自应单独处理，所以不能适用“部分实行、全部负责”原则。由于无法查明是谁过失致人死亡，根据存疑时有利于被告原则，每个人都不需要对死亡结果负责。如此，每个人都不构成故意伤害罪（过失致人死亡）。基于此，正确答案应当是B项。然而，官方答案认为B项说法错误。显然官方没有采取“要求故意且无罪说”。\footnote{要求故意且无罪说（部分犯罪共同说）中有种观点认为，打野猪案是孤立的过失致人死亡，所以甲乙都对死亡结果不用负责。而本题中，甲乙的过失致人死亡是在故意伤害罪的基础上派生的，所以对死亡结果都要负过失责任。然而，这种说法前后矛盾，不能实质性回答：为何孤立的过失致人死亡案要作无罪处理，而将过失致人死亡案镶嵌到故意伤害案中就作有罪处理？故意伤害罪（过失）致人死亡这种结果加重犯的原生态结构是，一个行为同时触犯两个罪，一是故意伤害罪，二是过失致人死亡罪，原本要想象竞合择一重罪论处，法律特事特办，按照结果加重犯处理。甲乙对故意伤害罪构成共同正犯，只意味着，针对故意伤害行为导致的伤害结果启动“部分实行、全部负责”原则，不意味着，针对过失行为导致的死亡结果也启动“部分实行、全部负责”原则。所以，要求故意且无罪说要一以贯之的话，对甲乙过失导致的死亡结果只能是都不负责。}

“要求故意且有罪说”认为，由于甲乙在“致人死亡”这个环节有意思联络，有相互的心理性贡献，所以无论是甲导致死亡结果，还是乙导致死亡结果，其中一方对另一方的致人死亡都有心理性贡献，所以对另一方导致的死亡结果都要负责。因此，甲乙都要对“过失致人死亡”负责，各自单独负责。在“致人死亡”环节，丙与甲乙有意思联络，对甲乙导致死亡结果有心理性贡献，所以也要对“过失致人死亡”负责。因此，三人构成故意伤害罪的共同犯罪，但各自单独构成结果加重犯“故意伤害罪（过失）致人死亡”。\CJKunderwave{官}\CJKunderwave{方采纳了这个结论。这表明官方采纳了“要求故意且有罪说”}

[总结]\CJKunderwave{无法查明的情形}

\InsertDiagram{figures/manual_table_p140_1.tex}

% Source PDF physical page 141; see audit/source_coverage.json.
\subsection*{三、共同正犯的种类}

上文讲解的共同正犯都是共同实行犯。共同正犯还有其他种类，亦即\CJKunderwave{其他起关键作用}\CJKunderwave{的共同正犯}，这主要是指，不是实行犯，但起到关键作用的犯罪人。

1.物理作用型共同正犯

例如，甲乙共同伤害丙，甲从后面死死抱住丙，乙拳击丙的头部，致丙重伤。乙的行为是故意伤害罪的实行行为，是实行犯。如何评价甲的作用？

甲的行为不是故意伤害罪的实行行为，不是实行犯。但是，甲的作用比一般的帮助犯的作用大。如果将甲认定为帮助犯，则量刑时只能按照从犯处罚，处罚会太轻。甲的作用大于帮助犯，小于实行犯，介于二者之间。可以将甲认定为共同正犯，如此，量刑时可以按照主犯处罚。

[巩固练习】题1，甲破解开保险柜的门，乙从里面拿钱。乙的行为是盗窃罪的实行行为（将财物转移占有），是实行犯。甲的行为不是盗窃罪的实行行为，但其行为起到了关键作用，因此构成共同正犯。

题2，甲将丙堵在胡同中，乙上前杀了丙。乙是故意杀人罪的共同实行犯，甲是共同正犯。

2.心理作用型共同正犯

（1）共谋共同正犯

例1，甲乙共谋盗窃丙家，计划一起实行。甲在去的路上出了车祸，未能去。乙顺利盗窃既遂。乙是实行犯，甲是共谋共同正犯，对乙的既遂结果要负责，也构成既遂，而不是犯罪预备。

例2，甲乙共谋盗窃丙家，由甲设计盗窃方案，由乙执行。乙盗窃既遂。甲构成共谋共同正犯，构成既遂。

（2）老大指派型共同正犯

例如，黑社会性质组织的老大指派小弟去盗窃丙家，小弟执行成功。老大的作用大于一般的教唆犯，因为有“老大”的组织权势的影响力，但影响力没有达到间接正犯的强制支配的程度，因此定共同正犯，由此量刑可按照主犯处罚。注意，这里的老大，只能是指犯罪组织的老大。如果只有甲乙两人犯罪，甲是大哥，指派乙犯罪，甲只能是教唆犯。

\InsertDiagram{figures/manual_p141_1.tex}

[提示】由于正犯主要是指实行犯，因此，除了特别说明，后文所称的共同正犯主要指实行犯。

% Source PDF physical page 142; see audit/source_coverage.json.
\section{间接正犯}

间接正犯与教唆犯的相同点在于，在客观要件上，都引起他人制造违法事实（A）。二者区别在于，间接正犯在此基础上，还对他人有支配力（B），他人被利用，成为工具人。教唆犯的客观要件是A，间接正犯的客观要件是A+B。

\InsertDiagram{figures/manual_table_p142_1.tex}

\subsection*{一、客观要件}

一个人对他人能形成支配力，主要源于三种情形：一是强制手段，二是欺骗手段，三是法定身份。

\subsubsection*{（一）强制手段}

1.迫使无责任年龄的人实施犯罪

例如，甲拐骗到一个8岁小孩后，迫使其去盗窃。甲除了构成拐骗儿童罪外，还构成盗窃罪的间接正犯。甲对小孩有支配力，小孩成为甲实施盗窃的工具。如果甲没有迫使，而是欺骗，则也构成间接正犯，属于通过欺骗而获得支配力。

2.强迫他人实施犯罪

例1，甲用枪指着乙，强迫乙猥亵丙。乙被迫照办。甲构成强制猥亵罪的间接正犯。乙可构成受强制的紧急避险。

例2，狱警强迫罪犯甲虐待罪犯乙，若甲不照办，就打断甲的腿。甲被迫照办。狱警构成虐待被监管人罪的间接正犯。甲可构成受强制的紧急避险。

3.强迫被害人实施自损行为

例如，君叫臣死，臣不得不死，大臣被迫自杀。皇上对大臣具有支配力，构成故意杀人罪的间接正犯。

\subsubsection*{（二）欺骗手段}

1.引诱无责任年龄、无责任能力的人

例如，甲引诱精神病患者乙窃取财物。甲构成盗窃罪的间接正犯。由于乙缺乏辨认能力，被甲当作犯罪工具。如果甲强迫乙去盗窃，也构成间接正犯，属于通过强制手段形成支配力。

2.欺骗他人，利用他人的不知情

（1）他人只有过失

例1（医生毒针案），医生将毒针交给护士，欺骗护士，吩咐其给病人注射，护士本应按规定检查针剂，但因为过于相信医生的权威而未检查并照办。医生是故意杀人罪的间接正犯，护士构成医疗事故罪。

例2，甲明知前方是王某，却对乙谎称前方有野兽，递给乙一把猎枪让乙射击。乙没有仔细察看便射击，打死了王某。乙构成过失致人死亡罪，甲构成故意杀人罪的间接正犯。

% Source PDF physical page 143; see audit/source_coverage.json.
（2）他人还有其他故意

行为模型：利用他人犯A罪的故意，实现自己犯B罪的间接正犯的目的。

例1（日本屏风案），乙不知道丙坐在高档屏风后面，而甲知道，甲为了杀死丙，唆使乙向屏风开枪，屏风被打碎，丙也中弹身亡。对甲、乙该如何处理？乙成立故意毁坏财物罪（实行犯）和过失致人死亡罪（假设有过失），想象竞合，择一重罪论处。甲成立故意毁坏财物罪（教唆犯）和故意杀人罪的间接正犯，想象竞合，择一重罪论处。甲利用乙犯故意毁坏财物罪的故意，实现了自已故意杀人罪的间接正犯的目的。

例2（2002年第38题），甲交给乙一包面粉，谎称是毒品，让乙卖给丙，所得货款一起分。乙不知情，竟然成功卖给了丙。丙不知情，支付了1万元。乙不构成贩卖毒品罪，属于不能犯。乙没有诈骗丙的钱财的故意，不构成诈骗罪。甲构成诈骗罪的间接正犯。甲利用乙犯贩卖毒品罪的故意，实现自己诈骗罪的间接正犯的目的。

例3（2019年试题），甲欲杀妻子丙，知道丙有血友病，隐瞒这一点，指使乙：“你将丙打成轻伤流血，我给你3万元。”乙将丙打成轻伤流血，无法预料丙流血致死。乙构成故意伤害罪，对死亡结果属于意外事件。甲构成故意杀人罪的间接正犯，利用乙犯故意伤害罪的故意，实现自己故意杀人罪的间接正犯的目的。

3.欺骗被害人实施自损行为

例1，甲带着毁坏财物的意图，欺骗乙：“你家那宠物狗已感染传染病，赶紧丢弃了！”乙相信便将价值1万元的宠物狗丢弃。甲构成故意毁坏财物罪的间接正犯。

例2，丈夫A和妻子B吵架后离家出走，有杀B故意的邻居C告诉B：“你假装上吊，我马上打电话叫A回来看看，吓吓他，让他以后不敢再和你争吵。”B听从C的意见，将搭在房梁上的绳子套在脖子上，C便离开，B很快吊死。C就属于利用被害人的行为达到杀人目的的间接正犯。

例3（2015年第56题），甲明知前方有陷阱，欲杀害乙，欺骗乙开车前往。乙掉入陷阱死亡。甲构成故意杀人罪的间接正犯。

\subsubsection*{（三）法定身份}

行为模型：有身份者利用无身份者，有身份者构成间接正犯。

例如（虐囚案），甲（监管人员）对乙（被监管人员）说：“你殴打虐待丙（被监管人员），我给你晚饭加个鸡腿。”乙照办，将丙打成轻伤。第一，从乙的角度看，就故意伤害罪而言，乙构成实行犯，甲构成教唆犯。第二，从甲的角度看。虐待被监管人罪是身份犯，身份是监管人员。按理说，由于甲对乙没有支配力，应该定该罪的教唆犯，但是，有教唆犯就应有实行犯，而乙因为没有监管人员的身份，不能构成该罪的实行犯。因此只能将甲扶上间接正犯的宝座。这种支配力源于法律规定，法律规定甲有特定身份。相应的，乙构成虐待被监管人罪的帮助犯，没有身份，不能构成实行犯，但能构成帮助犯。甲乙同时触犯两个罪，想象竟合，择一重罪论处。

[注意］间接正犯与被利用人的关系。

早期理论认为，间接正犯是一种单独犯罪的现象，间接正犯与被利用人不构成共同犯罪。这种看法过于绝对与片面，早已被槟弃。间接正犯也是正犯的一种，只要是正犯，就有资格给自己配备共犯（教唆犯、帮助犯）。例如，上述“虐囚案”中，甲构成虐待被监管人罪的间接正犯，乙构成虐待被监管人罪的帮助犯。二者构成该罪的共同犯罪。

% Source PDF physical page 144; see audit/source_coverage.json.
\subsection*{二、主观要件}

间接正犯，主观上必须是故意。这是因为，间接正犯是指使他人，利用支配他人犯罪，“指使、利用支配”这些动词本身只能是故意为之。因此，间接正犯的罪名只能是故意犯罪。

例如，医生疏忽大意，将一枚有毒针剂交给护士，让护士注射给病人。护士也疏忽大意，导致病人死亡。护士构成医疗事故罪。医生不可能构成故意杀人罪的间接正犯。医生也构成医疗事故罪，但不能说“医生构成医疗事故罪的间接正犯”。\footnote{提示：虽然在客观要件上，教唆犯与间接正犯是A与A+B的关系，但是整体上，二者并非一概是包容评价关系。（1）如果被利用人是故意制造违法事实，则间接正犯可以包容评价为教唆犯。例如，甲教唆8岁的乙盗窃，乙照办。乙是故意制造违法事实。甲构成盗窃罪的间接正犯，可以包容评价为盗窃罪的教唆犯。（2）如果被利用人是过失制造违法事实，则间接正犯难以包容评价为教唆犯。例如，医生欺骗护士，将毒针谎称正常针剂，让打给病人。护士疏忽大意打给病人。护士构成过失犯罪。医生构成故意杀人罪的间接正犯。该间接正犯难以包容评价为故意杀人罪的教唆犯。这是因为，教唆犯要有双重故意，也即故意教唆他人实施故意犯罪。对此，在后文教唆犯中会讲解。}

\section{正犯与共犯的关系}

这里的共犯是指狭义的共犯，也即教唆犯和帮助犯。在讲解教唆犯、帮助犯之前，需要先明确正犯与教唆犯、帮助犯的关系。

\subsection*{一、共犯从属性}

狭义共犯是教唆犯和帮助犯的统称，也简称共犯。关于共犯与正犯的关系，存在共犯独立性说与共犯从属性说。共犯从属性说是多数说。

例1，甲教唆乙盗窃，乙表面答应，实际上没有去。次日，乙欺骗甲盗窃了。

例2，甲帮助乙盗窃，在墙外望风，实际上乙没有实施盗窃，而是回家睡觉了。次日，乙欺骗甲盗窃了。

共犯独立性说认为，对教唆者、帮助者的定罪应独立判断，不用考虑实行者是否构成犯罪。甲实施了教唆行为，便构成教唆犯，属于犯罪未遂。甲实施了帮助行为，便构成帮助犯，属于犯罪未遂。

共犯从属性说认为，实行者是直接实施犯罪的人，对法益的侵害具有直接性；教唆者、帮助者只是促使犯罪、帮助犯罪的人，对法益的侵害具有间接性，需借助实行者去实施犯罪才能实现对法益的侵害。如果实行行为缺位，教唆行为、帮助行为便失去了凭借，对法益便不具有侵害的危险。这样，教唆行为、帮助行为便不会成为危害行为，也就不构成犯罪。结论：如果实行者没有法益侵害性（违法性），不构成犯罪，那么教唆者、帮助者也没有法益侵害性，也不构成犯罪。实行者构成犯罪是教唆者、帮助者构成犯罪的前提条件（必要条件）。也即，在犯罪成立条件上，正犯具有主导地位，共犯（教唆者、帮助者）从属于正犯，处于从属地位。共犯从属性说是法考的官方立场。（2009年主观题）

\subsection*{二、犯罪形态的从属性}

上述从属性是就犯罪成立条件而言的。在犯罪形态上，共犯也具有从属性，也即，教唆犯、帮助犯的犯罪形态取决于实行犯。

% Source PDF physical page 145; see audit/source_coverage.json.
\CJKunderwave{结论}：（1）当实行犯尚在预备阶段，则教唆犯、帮助犯只能在预备阶段确定犯罪形态；（2）当一个实行犯进入实行阶段，则教唆犯、帮助犯、其他实行犯只能在实行阶段确定犯罪形态。\footnote{这是一般而言。帮助犯要在实行阶段找犯罪形态，要求帮助作用维持到实行阶段。对此，帮助犯一节有讲解。}

［注意］犯罪中止具有独立性，也即某个共犯人构成犯罪中止，其效力不能惠及于其他共犯人。这是因为，对中止犯之所以从宽处罚，是因为其人身危险性降低了，而人身危险性是人身专属性特征。某个共犯人的人身危险性降低，并不意味着其他共犯人人身危险性降低。所以，共犯人要成立犯罪中止，要求亲力亲为。关于共犯人的犯罪中止条件，参见后文“中途退出：共犯关系的脱离”。

\InsertDiagram{figures/manual_table_p145_1.tex}

例1（2003年第42题），乙要杀人，让李某提供毒药，李某提供了，乙又放弃投毒杀人，将毒药扔掉。乙构成预备阶段的犯罪中止。乙的犯罪中止的效力只能及于乙自已。乙的犯罪中止，对于李某而言是意志以外原因，李某构成犯罪预备。

例2，甲入户盗窃，乙蹲守望风。乙在外面被抓，甲不知道乙被抓，在里面盗窃既遂。乙对甲的既遂有心理性贡献，构成既遂。大致而言，如果有一个人既遂，大伙就都既遂了。\footnote{这只是大致结论。如果教唆犯、帮助犯没有引起、促进正犯制造既遂结果，则不构成教唆犯既遂、帮助犯既遂。对此，后文有详细讲述。}其他人就别想再成立中止或未遂。当然，乙如果在甲既遂前，主动消除自已的贡献，脱离共犯关系，则成立犯罪中止。对此，后文有讲解。

[巩固练习］题1，甲乙合谋抢劫丙，使用暴力后，发现丙身上没带钱，甲乙便押着丙前往丙家中取钱。途中，甲因临时有事要离开，便让乙继续前往丙家。乙到了丙家，发现丙生活困难，于心不忍，遂放弃。乙构成实行阶段的犯罪中止。乙的中止效力不能及于甲。乙的中止，对于甲而言属于甲意志以外的原因。甲因为意志以外原因而未能得逞，构成犯罪未遂。（2022年试题）

题2，甲乙欲共同抢劫丙的财物，在去往丙家途中，甲忽然肚子疼，便让乙继续前往。乙进入丙家中，进行抢劫，丙哀求，乙便放弃抢劫。乙构成实行阶段的犯罪中止。在预备阶段，甲虽然未能着手实行，但是当共同正犯乙着手实行后，甲的犯罪形态便只能在实行

% Source PDF physical page 146; see audit/source_coverage.json.
阶段找。乙的中止效力不能及于甲。乙的中止，对于甲而言属于甲意志以外的原因。甲因为意志以外原因而未能得逞，构成犯罪未遂。（2022年试题）

题3，甲乙计划入户抢劫丙，甲是实行者，乙是帮助者。甲提前一天来丙家踩点，发现丙是自己中学同学，便心生悔意，将和乙一起抢劫丙的计划向丙如实交代。丙欲抓住乙，便告知甲继续进行“犯罪计划”。次日，甲乙来到丙家，乙在外望风。甲进入丙家后和丙一起报警。警察赶到将望风的乙抓获。甲构成预备阶段的犯罪中止。这是因为，甲虽然没有向乙明确告知退出意思，但是，甲实施了抓捕乙的行为。该行为比“明确告知”更符合犯罪中止的实质要求。犯罪中止的实质要求是降低人身危险性。乙作为帮助犯构成犯罪预备。（2024年试题）

\section{教唆犯}

[问题］甲教唆乙杀害丙，乙答应，来到丙家小区，由于样子鬼鬼崇崇，被门卫老大爷质问，乙为了摆脱纠缠，一气之下杀死老大爷，并逃离。对甲、乙该如何处理？（杀老大爷案)

教唆犯，是指故意引起他人制造违法事实（法益侵害事实）。例如，甲唆使乙杀人，乙答应，实施了杀人。乙是直接实行犯，甲是教唆犯。教唆犯的关键词是“引起”。

\InsertDiagram{figures/manual_table_p146_1.tex}

\subsection*{一、客观要件}

在客观要件上，要求教唆行为引起正犯制造违法事实（法益侵害事实），这包括两个层次：第一，\CJKunderwave{教唆行为引起正犯的违法行为。这是教唆犯的成立条件}。第二，\CJKunderwave{教唆行为引}\CJKunderwave{起正犯的违法结果。这是教唆犯的既遂条件}。重点是判断引起与被引起的关系（因果性）。

\subsubsection*{（一）教唆犯的成立条件：教唆行为引起正犯的违法行为}

1.他人已经具有犯某罪的意图，此时教唆其犯该罪，不构成教唆犯，因为二者不存在引起与被引起的关系。例如（2013年第55题），甲教唆乙杀害小芳，乙表面答应，实际上乙早有此意，乙杀害了小芳。甲不成立教唆犯，因为乙的杀害行为不是甲的教唆行为引起的。

2.他人已经有犯轻罪的意图，教唆他人犯重罪，则教唆者构成重罪的教唆犯。例如，乙本想诈骗，甲教唆乙：“诈骗多累啊，直接抢劫最快！”乙便实施抢劫。甲构成抢劫罪的教唆犯。

3.他人已有犯基本犯的意图，教唆他人犯情节加重犯，则教唆者构成情节加重犯的教唆犯。例如，乙本想抢劫，甲教唆乙持枪抢劫，乙便持枪抢劫。甲构成持枪抢劫的教唆犯，也即对乙适用持枪抢劫的法定刑，对甲也适用持枪抢劫的法定刑。这是因为，非法持枪已经是一个不同的犯罪情节。

4.他人已有犯基本犯的意图，教唆他人犯数额加重犯，则教唆者构成数额加重犯的

% Source PDF physical page 147; see audit/source_coverage.json.
（心理性）帮助犯。例如，乙本想盗窃数额较大的财物，甲教唆乙盗窃数额巨大的财物，乙便照办。甲构成盗窃数额巨大财物的帮助犯。这是因为，甲没有教唆乙实施新的犯罪行为或情节，只是提高了犯罪数额，只是一种心理性的促进作用。

5.他人已经有犯重罪的意图，教唆他人犯性质相同的轻罪，则教唆者不构成教唆犯，也不构成心理性帮助犯，因为降低了法益受到的危险。例如，乙本想抢劫，甲教唆乙：“抢劫要判死刑的，还是偷吧！”乙便实施盗窃。甲不构成犯罪。

他人已经有犯重罪的意图，教唆他人犯性质不同的轻罪，则教唆者构成该轻罪的教唆犯。例如，乙本想抢劫，甲教唆乙：“别去抢劫，还是去强制猥亵吧！”乙便强制猥亵。甲构成强制猥亵罪的教唆犯。

\subsubsection*{（二）教唆犯的既遂条件：教唆行为引起正犯的违法结果}

例如，本节开头“杀老大爷案”中，（1）先分析乙：乙对预定目标丙构成故意杀人罪（犯罪预备），对实害对象（老大爷）构成故意杀人罪（既遂），整体是一个行为（杀老大爷也是为了杀丙计划顺利进行），想象竞合，择一重罪论处，定故意杀人罪（既遂）。

（2）后分析甲：对预定目标丙，由于乙构成故意杀人罪（犯罪预备），所以甲也构成故意杀人罪（犯罪预备）。对实害对象（老大爷），由于杀死老大爷是乙自己节外生枝，与甲没有关系，甲不用负责。结论：甲的教唆行为引起了乙的违法行为（杀丙的行为），因此甲成立教唆犯；但是，甲的教唆行为并没有引起乙的违法结果（杀死老大爷），因此甲不会构成故意杀人罪既遂。

[提示】传统理论往往认为：“正犯既遂，则教唆犯一定既遂；一人既遂，则全部既遂。”这种说法泛泛而论也是正确的，但是具体分析，则并不准确。教唆犯要构成既遂，是有条件的，也即教唆行为引起正犯的违法结果。

\subsection*{二、主观要件}

在主观要件上，要求教唆犯必须具有教唆故意。教唆故意，是指\CJKunderwave{故意引起}他人实施故意犯罪。这里有两个故意，一是，实行犯是故意犯罪，二是，教唆犯是故意教唆。缺少一个故意，教唆者不构成教唆犯。实际上，“教唆、唆使”这些动词本身就意味着故意为之。教唆犯承受的罪名只能是故意犯罪。

\subsubsection*{（一）不是故意引起，不构成教唆犯}

例1，甲和乙聊天时，说道：“丙强奸了我的女神小芳，我一定要杀死丙！”由于小芳也是乙的女神，乙便暗下决心要杀死丙，并且在甲动手前杀死了丙。在客观上，甲引起了乙制造违法事实。但是主观上，甲不是故意引起乙制造违法事实，并且，甲对乙的行为及其结果，无法预见，因此甲不构成犯罪。

例2（女友回家案），女友看到男友已经喝醉，仍执意要求男友驾车送其回家。男友因醉驾不慎撞死人。男友构成交通肇事罪。女友也构成交通肇事罪，理由是，第一，客观上，女友指使醉酒者驾驶，由此给道路公共安全制造危险。这种指使行为是一种违法行为。第二，主观上，女友意识到醉酒者驾驶有可能撞死人，不过女友对这种死亡结果的发生不是故意心理，但是存在过失心理，属于过于自信的过失。概言之，女友对男友的驾驶行为存在指挥过失，应对男友制造的违法结果负责。

这表明，\CJKunderwave{过失犯罪中，对结果负责有两种情形：一是对自己亲自制造的结果负责，这}

% Source PDF physical page 148; see audit/source_coverage.json.
\CJKunderwave{是一种事实判断。二是对他人制造的结果负责}，其条件是，自己与他人有意思联络，并对他人的违法结果存在过失行为，表明二人一起制造违法事实。这是一种法律评价。这也表明，交通肇事罪不仅可以由司机构成，也可以由其他人构成。注意，不能说“女友构成交通肇事罪的教唆犯”，因为交通肇事罪是过失犯罪。女友构成交通肇事罪的单独正犯。（2019年试题）

\subsubsection*{（二）故意教唆实施不能犯，不构成教唆犯}

故意教唆他人实施不能犯（没有法益侵害危险的行为），不构成教唆犯。

例1，甲为了试探乙的胆量，给乙一把空枪，乙不知是空枪。甲让乙开枪打死不远处的丙，乙便开枪。

例2，甲给了乙一把有子弹的枪，让乙杀死丙。乙开枪，但子弹卡壳。

例1中，因为乙的行为没有任何危险，乙无罪，根据共犯从属性，甲不构成教唆犯，也无罪。乙是不能犯，甲也是不能犯。不能犯是指主观上有犯罪故意，但客观上没有法益侵害的危险。理论上将故意教唆他人实施不能犯这种情形也称为未遂的教唆。

例2中，乙的行为有法益侵害的危险，作为实行犯，构成故意杀人罪未遂。甲作为教唆犯，也构成故意杀人罪未遂。对甲这种犯罪情形的称呼，理论上的称谓有“教唆犯未遂”“未遂的教唆犯”。

\subsection*{三、处罚}

1.第29条第1款第1句：“教唆他人犯罪的，应当按照他在共同犯罪中所起的作用处罚。”

（1）教唆犯起主要作用，定主犯；起次要作用，定从犯。教唆犯也有可能成为胁从犯。例如，甲掌握了乙（妇女）的裸照，威胁乙去教唆丙（法官）徇私枉法。乙被迫答应。乙既属于胁从犯，又属于教唆犯。\footnote{提示：甲雇佣乙杀丙，乙杀了丙。雇佣行为属于教唆行为，定教唆犯。（1）对雇佣者不能定共同正犯，因为雇佣者的威力达不到黑社会老大指派小弟的程度。（2）对雇佣者不能定间接正犯，因为雇佣者对受雇者没有支配力、控制力。（3）雇佣者这种教唆犯，比一般教唆犯所起作用更大一些，量刑时可以定主犯。}

（2）教唆犯可以不止一人。例如，甲、乙共同教唆丙盗窃，丙便盗窃。甲、乙均构成教唆犯，如果甲起到主要作用，则构成主犯。

（3）教唆对象必须特定，如果不特定，就是“煽动”。例如，刑法第373条规定的煽动军人逃离部队罪，对象不能特定；如果特定，就是第435条规定的逃离部队罪的教唆犯。

（4）间接教唆的，构成教唆犯。例如，甲教唆乙，让乙教唆丙去盗窃，乙便教唆了丙，丙实施了盗窃。甲属于间接教唆，构成盗窃罪的教唆犯。

2.第29条第1款第2句：“教唆不满十八周岁的人犯罪的，应当从重处罚。”

该规定包括三种情形：

（1）甲教唆乙（17周岁）盗窃，乙照办。甲构成教唆犯。

（2）甲教唆乙（15周岁，但证明有规范意识、独立作案能力）盗窃，乙照办。甲构成教唆犯。甲对乙没有支配力，不构成间接正犯。

% Source PDF physical page 149; see audit/source_coverage.json.
（3）甲教唆乙（8周岁，推定没有规范意识和独立作案能力）盗窃，乙照办。甲构成教唆犯。并且，甲对乙具有支配力，构成间接正犯。（2013年第9题）

3.第29条第2款：“如果被教唆的人没有犯被教唆的罪，对于教唆犯，可以从轻或者减轻处罚。”

对该款的理解，共犯独立性理论认为：即使被教唆者没有犯罪，教唆者也构成犯罪，只是可以从轻或减轻处罚。这是少数说。根据共犯从属性理论（多数说），该款是指被教唆的人（实行犯）已经着手实行了犯罪，但是没有达到既遂的程度，只是未遂或中止了，所以可以从宽处罚。

例如（2009年主观题），乙教唆甲抢劫甲的情妇丙的钱，甲说：“别瞎整！”乙未再声。过几天，甲去丙家做客，临时起意抢劫丙的钱。第一，甲拒绝接受教唆。乙不构成教唆犯，不能给乙定教唆犯未遂。第二，甲几天后的抢劫罪是单独犯罪，与乙无关。不能因为甲构成抢劫罪既遂，就给乙定抢劫罪教唆犯既遂。

\OriginalPanel{典型真题}

\ensuremath{>}\ensuremath{>}

甲与乙女有婚外情，甲欲与妻子丙离婚。乙唆使甲用毒牛奶杀害丙。甲接受了教唆。几天后，甲将一瓶毒牛奶递给丙。丙不知道牛奶有毒，又将牛奶递给身边的孩子丁喝。甲见状，就说了一句“他喝过了，不用喝了”便走开了，没有实施其他阻止行为。丁喝了毒牛奶后死亡。下列说法错误的是？（2021年回忆版金题）\footnote{[答案]甲对妻子丙构成故意杀人罪未遂，对孩子丁构成不作为的故意杀人罪既遂，这是前后两个独立的行为，因此应当并罚。当实行犯甲对丙构成故意杀人罪未遂，那么教唆犯乙对丙也构成故意杀人罪未遂。乙的教唆行为与正犯甲的违法结果（丁的死亡）之间缺乏引起与被引起的因果关系，所以不能将该死亡结果归责给乙，乙对丁的死亡不承担教唆犯的责任，不能因为丁的死亡而给乙定故意杀人罪既遂。本题答案为D。}

B.甲对丁构成故意杀人罪既遂A.甲对丙构成故意杀人罪未遂

C.乙对丙构成故意杀人罪未遂D.乙对丁构成故意杀人罪既遂

\section{帮助犯}

[问题］甲欲盗窃丙家，让乙提供有用的钥匙，乙将有用的钥匙交给了甲。甲拿着钥匙开门，由于操作不当，一时打不开，情急之下破门而入，盗窃既遂。对甲乙该如何处理？(提供钥匙案）

帮助犯，是指故意促进他人制造违法事实。例如，乙要杀人，甲为其提供凶器。乙拿着甲的凶器杀了人。乙是直接实行犯，甲是帮助犯。帮助犯的关键词是“促进”。

\InsertDiagram{figures/manual_table_p149_1.tex}

% Source PDF physical page 150; see audit/source_coverage.json.
\subsection*{一、客观要件}

\subsubsection*{（一）帮助犯的成立条件：帮助行为促进正犯实施违法行为}

成立帮助犯，帮助行为需要具备两个条件：第一，帮助行为本身具有可能的促进作用；第二，这种合格的帮助行为连接到（作用于）正犯的违法行为上。

1.帮助行为具有可能的促进作用，是指帮助行为本身具有法益侵害的危险性、可能性。

例如，甲绑架了乙的智障孩子丙，然后让保姆照料丙，给丙做饭洗衣，并告知保姆绑架事实。保姆不构成绑架罪的帮助犯，因为照料行为有益于丙，没有法益侵害的危险性。

模型1，甲欲盗窃丙家，让乙提供有用的钥匙，乙故意提供一把根本无法用的钥匙。甲使用之，无效。乙不构成帮助犯。又如，乙想帮助，找配钥匙师傅，师傅欺骗乙，故意配一把根本无法用的钥匙。乙不知情，提供给甲。甲使用之，无效。乙不构成帮助犯。

2.合格的帮助行为连接到（作用于）正犯行为。

模型2，甲欲盗窃丙家，让乙提供有用的钥匙。乙拿着有用的钥匙去甲家。情形一：在路上乙的钥匙被小偷偷去了。甲自已盗窃既遂。情形二：乙到了甲家，甲不在家，甲打电话让乙将钥匙放在自已的信箱里。乙不慎将钥匙放到甲邻居的信箱里。甲未拿到钥匙，自己盗窃既遂。（2014年第54题）

这两种情形中，乙的帮助行为本身具有法益侵害的危险性、可能性，是合格的帮助行为，但由于意志以外原因未能连接到（作用于）正犯行为的危险流中，便不可能发挥促进作用，不可能对法益（丙家财物）具有危险性，因此不成立帮助犯。这种情形属于帮助行为本身因意志以外原因未能成功、未得逞。理论上将这种现象称为\CJKunderwave{帮助未遂}。\footnote{2014年第54题C项：“乙欲盗窃汽车，让甲将用于盗窃汽车的钥匙放在乙的信箱。甲同意，但错将钥匙放入丙的信箱，后乙用其他方法将车盗走”。官方答案认为，甲构成盗窃罪帮助犯未遂。该答案值得商。甲的帮助行为未能连接到正犯行为的危险流中，属于帮助未遂，不成立帮助犯，不成立犯罪，而非犯罪未遂。参见张明楷：《共同犯罪的认定方法》，载《法学研究》2014年第3期。}

\subsubsection*{（二）帮助犯的既遂条件：帮助行为促进正犯制造违法结果（因果性）}

当帮助行为连接到正犯行为上，帮助行为便成立帮助犯。当帮助行为与正犯结果具有因果性时，帮助犯便既遂。也即，在正犯行为的危险流导致结果的过程中，帮助行为要发挥实际贡献，帮助犯才既遂。

模型3，甲欲盗窃丙家，让乙提供有用的钥匙。乙将有用的钥匙给了甲。甲在去丙家途中，钥匙被盗，用别的办法盗窃既遂。乙的帮助行为连接到甲的正犯行为，成立帮助犯。但该帮助行为对正犯行为发挥的作用只是维持到了正犯行为的预备阶段，便没能继续发挥作用，因此，该帮助行为只能构成犯罪预备。对于帮助犯的犯罪预备，实务中一般不予定罪处罚。

[练习］甲欲盗窃丙家，让乙提供有用的钥匙。乙将有用的钥匙给了甲。甲拿着该钥匙来到丙家，看到丙家门开着，屋里没人，甲直接走进去，盗窃既遂。乙的帮助行为连接到甲的正犯行为，成立帮助犯。但该帮助行为对正犯行为发挥的作用只是维持到了正犯行为的预备阶段，没能维持到着手实行阶段，因此，该帮助行为只能构成犯罪预备。

% Source PDF physical page 151; see audit/source_coverage.json.
模型4，本节开头“提供钥匙案”中，甲欲盗窃丙家，让乙提供有用的钥匙，乙将有用的钥匙交给了甲。甲拿着钥匙开门，由于操作不当，一时打不开，情急之下破门而入，盗窃既遂。乙的帮助行为连接到甲的正犯行为，成立帮助犯。并且，该帮助行为对正犯行为发挥的作用维持到了正犯行为的着手实行阶段。但是，在正犯行为的危险流导致结果的过程中，乙的帮助行为没有发挥实际贡献，因此构成犯罪未遂，而非既遂。（2013年第54题，2017年第6题）

模型5，甲欲盗窃丙家，让乙提供有用的钥匙。乙将有用的钥匙给了甲。甲拿着该钥匙来到丙家，用钥匙打开门，进去找到财物，盗窃既遂。乙的帮助行为对正犯行为发挥的作用维持到着手实行阶段，并且对甲的既遂结果发生了实际贡献，构成既遂。

\InsertDiagram{figures/manual_p151_1.tex}

[提示］称谓问题。“帮助犯未遂”与“未遂的帮助犯”是针对同一案件事实的不同称谓。“教唆犯未遂”与“未遂的教唆犯”也是针对同一案件事实的不同称谓。\hypertarget{extended-note-151-1-ref}{}\textsuperscript{\hyperlink{extended-note-151-1}{附注}}

\subsubsection*{（三）心理性帮助}

帮助行为的方式有两种。第一，物理性帮助。例如，提供工具。第二，心理性帮助，这是指对正犯产生心理性促进作用。

\par\noindent\hypertarget{extended-note-151-1}{\hyperlink{extended-note-151-1-ref}{\textbf{附注}}}\quad 在此应注意名与实的关系。先明确案件事实，再看对该事实的称谓是什么。对同一个事实可以从不同侧面、不同角度赋予不同称谓。

\InsertDiagram{figures/manual_table_p151_2.tex}

% Source PDF physical page 152; see audit/source_coverage.json.
例1，甲得知乙欲杀丙，对乙讲：“你若真杀了丙，我奖励你5万元！”乙便很快杀了丙。甲构成故意杀人罪的帮助犯既遂。

例2（望风行为），甲、乙共谋盗窃，甲在外望风，乙入户盗窃。甲一方面提供心理性帮助，使乙在里面安心盗窃；另一方面也提供物理性帮助，比如主人回家，甲借故拖延主人进屋。

\InsertDiagram{figures/manual_table_p152_1.tex}

\subsection*{二、主观要件}

在主观要件上，要求帮助犯必须具有帮助故意。帮助故意，是指\CJKunderwave{故意促进}他人实施故意犯罪。这里有两个故意，一是，实行犯是故意犯罪，二是，帮助犯是故意帮助。缺少一个故意，帮助者不构成帮助犯。

例1，甲欲杀死丙，向猎人乙借枪，谎称去打猎。甲在乙心里一向是善良的人，乙便信以为真，将枪借给甲。甲用枪杀死丙。乙没有帮助甲杀人的故意，且无法预见甲的杀人行为，不构成犯罪。

例2（2017年第7题），甲欲杀死丙，对乙谎称：“我去盗窃丙家，你帮我望风。”乙不知情，答应照办。甲入户杀死丙。客观上，乙的行为给甲的杀人提供了帮助，但是主观上，乙没有帮助甲杀人的故意，不构成故意杀人罪的帮助犯。虽然乙有帮助甲盗窃的故意，但是客观上甲没有盗窃行为，乙也不构成盗窃罪的帮助犯。乙构成非法侵入住宅罪的帮助犯。

例3，甲乙共谋入户抢劫，乙望风，甲入户后发现户内没人，拿走财物，出来欺骗乙：“我浴血奋战，才抢到这点财物，给我多分点。”乙不知情而答应。客观上，乙的行为给甲的盗窃提供了帮助，但是主观上，乙只有帮助甲抢劫的故意，不过，帮助抢劫的故意可以包容评价为帮助盗窃的故意。因此，乙构成盗窃罪的帮助犯既遂。

% Source PDF physical page 153; see audit/source_coverage.json.
[总结］四大角色的主观要件

\InsertDiagram{figures/manual_table_p153_1.tex}

\subsection*{三、中立的帮助行为}

例如，出租车司机甲明确得知乘客乙要抢劫前方行人丙，仍将其运到丙的身边。乙下车抢劫了丙。甲的行为是属于民事义务行为，还是已经构成抢劫罪的帮助犯？

\CJKunderwave{判断标准}：第一，主观上是否\CJKunderwave{明确知道}对方在犯罪；第二，客观行为给对方犯罪是否起到了实质且紧迫的促进作用。基于此，甲构成抢劫罪的帮助犯。

例1，甲每天给乙开设的赌场送盒饭，不构成开设赌场罪的帮助犯，因为没有实质性帮助。

例2，甲来到商店买老鼠药，商店老板乙猜出甲可能用来毒杀妻子，仍予以出售。两天后，甲用老鼠药毒杀了妻子。乙不构成帮助犯，因为其一，乙不是明确知道，只是猜测；其二，乙的帮助没有紧迫的促进作用。

[巩固练习］题1（2013年第55题），乙、丙在五金店门前互殴，店员甲旁观。乙边打边掏钱向甲买一羊角锤。甲递锤时对乙说“你打伤人可与我无关”。乙用该锤将丙打成重伤。甲构成故意伤害罪的帮助犯。

题2，甲欠乙五万元，甲得知，若还钱，乙会用于购买毒品进行贩毒，仍还钱。次日，乙用于购买毒品。甲不构成贩卖毒品罪的帮助犯。第一，甲有还钱的义务，乙有讨债的请求权。还债行为具有中立性。第二，甲的行为没有起到紧迫的帮助作用。同理，甲欠乙五万元。后乙犯了罪，要求甲还钱。甲知道乙犯了罪，要潜逃，需要钱，仍向乙还钱。虽然甲的还钱在客观上帮助了乙逃匿，但甲不构成窝藏罪。

题3，甲出售专门用于窃电的设备，购买者用来窃电。该设备在技术上不具有中立性，而是专用性。甲构成盗窃罪的帮助犯。

% Source PDF physical page 154; see audit/source_coverage.json.
\section{正犯的责任年龄和责任能力}

前文已述，共同犯罪是指一起制造违法事实；到了排除责任（谴责）阶段，谴责谁，不谴责谁，独立评价。因此，就正犯（主要是实行犯）的成立条件而言，只要求其通过实行行为制造违法事实即可，至于其是否具有责任年龄、责任能力，则在所不问。

\subsection*{一、共同正犯}

1.正犯甲（有责任年龄）+正犯乙（无责任年龄）

例如，狗剩（20岁）和狗蛋（10岁）共同轮奸了小芳。根据中国化改良版的四要件体系，在“制造违法事实”阶段，违法具有一起性。二人构成该阶段的共同正犯。在排除责任（谴责）阶段，无法谴责狗蛋，狗蛋最终无罪。狗剩构成强奸罪的共同实行犯，在量刑上，按照“轮奸”情节加重处罚。若按照苏联四要件体系，狗蛋一开始就无罪，只能单独处理狗剩。如此对狗剩不能按照“轮奸”加重处罚，因为轮奸是指强奸罪的共同犯罪。这种结论显然不妥当。

2.正犯甲（有责任能力）+正犯乙（无责任能力）

例如，甲邀请乙一起伤害小芳，乙答应照办。实际上乙是精神病患者，甲不知道。小芳遭受重伤，能够查明，是某一脚将小芳脾脏踹破裂，但无法查明是甲踹的，还是乙踹的。根据中国化改良版四要件体系，在“制造违法事实”阶段，违法具有一起性，甲乙构成该阶段的共同正犯。根据“部分实行、全部负责”原则，甲乙都对重伤结果负责。在排除责任（谴责）阶段，无法谴责乙，乙最终无罪。甲最终定故意伤害罪（重伤）既遂。若按照苏联四要件体系，乙不是合格的犯罪主体，直接无罪。只能单独处理甲，那么根据存疑时有利于被告原则，甲对重伤结果不负责。这显然不合理。

\subsection*{二、帮助犯+正犯}

1.帮助犯（有责任年龄）+正犯（无责任年龄）

例如，10岁的狗蛋杀害小芳，指使老大爷为其望风。后狗蛋得逞。二人构成故意杀人罪的共同犯罪，狗蛋是正犯，老大爷是帮助犯。在排除责任（谴责）阶段，狗蛋未达到责任年龄，最终不负刑事责任，最终无罪。

2.帮助犯（有责任能力）+正犯（无责任能力）

例如，乙精神病发作，要打丙。甲故意给乙木棒，说：“用这个打！”乙接过木棒打丙。乙与甲构成故意伤害罪的共同犯罪，乙是正犯，甲是帮助犯。在排除责任（谴责）阶段，乙因为无责任能力而最终无罪；甲最终定故意伤害罪帮助犯。（2018年试题）

\subsection*{三、教唆犯+正犯}

\subsubsection*{（一）教唆犯(有责任年龄）+正犯（责任年龄分情况)}

1.实行者已达责任年龄，具有规范意识、独立作案能力。

例如，甲教唆乙（16周岁）盗窃，乙照办。甲对乙缺乏支配力，仅构成教唆犯(A)。

% Source PDF physical page 155; see audit/source_coverage.json.
2.实行者未达责任年龄，不具有规范意识、独立作案能力。

例如，甲教唆乙（8周岁）盗窃，乙照办。第一，在制造违法事实阶段，乙是“实行犯”，甲是教唆犯。第二，在排除责任（谴责）阶段，乙因缺乏责任年龄而无罪，甲负教唆犯的责任。第三，甲在教唆犯（A）的基础上，对乙具有支配力（B），因此构成间接正犯（A+B）。就高不就低，对甲最终以间接正犯论处。

3.实行者未达责任年龄，但证明具有规范意识（知道自已行为的法律后果）、独立作案能力，则意味着教唆者对实行者难以形成支配力，不构成间接正犯，而只构成教唆犯。

例如（15周岁盗窃案），乙15周岁，曾因盗窃被公安机关管教过，又经常盗窃。甲唆使乙盗窃，乙照办。此时，甲不构成间接正犯，而仅仅构成教唆犯。实行犯是乙，制造违法事实阶段的“实行犯”。乙最终无罪。

【提示］传统理论认为，教唆犯与间接正犯的区分标准是，看被利用人是否具有责任年龄。具有责任年龄，则指使者构成教唆犯；不具有责任年龄，则指使者构成间接正犯。也即，指使无责任年龄者犯罪，一律构成间接正犯。基于此，上述“15周岁盗窃案”中，甲不构成教唆犯，而构成间接正犯。这种观点的错误在于，第一，年龄只是个形式标准，应当采取实质标准，也即看有无支配力。第二，将教唆犯与间接正犯的关系理解成对立排斥关系。实际上，二者有时候是A与A+B的包容评价关系（位阶关系）。而且传统理论会导致定罪障碍。

例如（虚报年龄案），甲为儿子乙（15周岁）虚报年龄，考取公务员，成为一名警察。甲唆使乙刑讯逼供犯人丙。乙照办。（1）先分析乙，在制造违法事实阶段，乙是刑讯逼供罪的“实行犯”。在排除责任事由阶段，乙未达16周岁，最终无罪。（2）后分析甲。若依传统理论，甲不构成刑讯逼供罪的教唆犯，而构成刑讯逼供罪的间接正犯。但是，甲无司法工作人员身份，不能构成刑讯逼供罪的间接正犯。如此，对甲便无法处理。而根据实质标准，甲构成教唆犯。实行犯是乙，制造违法事实阶段的“实行犯”。

［总结］小孩犯罪，大人角色。

（1）10岁的狗蛋要盗窃，指使老大爷为其望风。老大爷照办。老大爷构成作为的帮助犯。

（2）10岁的狗蛋是好孩子，老大爷指使狗蛋去盗窃。狗蛋照办。老大爷构成作为的间接正犯。

（3）10岁的狗蛋在盗窃，父亲发现，能够轻易有效阻止，却不阻止。狗蛋是“制造违法事实”意义上的“实行犯（正犯）”。父亲构成不作为的帮助犯，并且在帮助犯的基础上构成共同正犯，因为对结果的发生所起的作用超过了帮助犯的角色，达到了共同正犯的程度。注意：多数说认为，父亲不构成间接正犯，因为父亲并没有指使狗蛋、支配狗蛋去犯罪，只是单纯地不阻止。（2021年试题）

\subsubsection*{（二）教唆犯（有责任能力）+正犯（无责任能力）}

例如（教唆伤害案），甲教唆乙伤害丙，实际乙是精神病患者，但甲对此不知情。在“制造违法事实阶段”，乙与甲构成故意伤害罪的共同犯罪，乙是正犯。甲构成教唆犯，但不构成间接正犯，因为甲没有实施间接正犯的故意。在排除责任（谴责）阶段，乙因为无责任能力而最终无罪；甲最终定故意伤害罪的教唆犯。

[注意】传统理论认为，具有责任能力者与无责任能力者共同犯罪，具有责任能力者

% Source PDF physical page 156; see audit/source_coverage.json.
一律构成间接正犯。这种看法是不正确的。上述“教唆伤害案”就印证了这一点。真题专门批判过这种看法。（2014年第10题）\footnote{“正犯+正犯”中，正犯也不需要具有责任年龄、责任能力。}

\section{共同犯罪的参与时间}

\subsection*{一、中途参与：承继的共同犯罪}

\subsubsection*{（一）参与时间}

承继的共同犯罪，是指中途参与他人犯罪。与中途参与不同的是，事前参与与事后参与。

\InsertDiagram{figures/manual_table_p156_1.tex}
\SourceFootnote{\textbf{表注②}\quad 有司法解释规定，抢劫行为导致被害人轻伤，即使没有取得财物，也构成抢劫罪既遂。该规定是特殊规定，不能推广适用，不适用于承继共犯场合。详细分析见分则“抢劫罪”一节。}
\SourceFootnote{\textbf{表注③}\quad 提示：（1）掩饰、隐瞒犯罪所得罪只能在上游犯罪既遂、产生犯罪所得之后实施（多数说）。（2）窝藏、包庇罪既可以在上游犯罪既遂后实施，也可以在上游犯罪构成其他终局形态（预备、中止、未遂）后实施。（3）帮助毁灭、伪造证据罪既可以在上游犯罪既遂后实施，也可以在上游犯罪实施期间实施，因为只要发生犯罪，就有了犯罪证据，不是等到了犯罪终局后才产生犯罪证据。}

\subsubsection*{（二）责任划分}

[结论】中途参与者只需对参与后的事情负责，对参与前的事情不负责。

例1，甲抢劫丙，将丙打倒在地，乙此时参与进来帮甲捡起了财物，丙经抢救无效死亡。事后查明，丙的死亡是由甲的暴力造成的。甲乙构成抢劫罪承继的共同犯罪，甲构成抢劫罪致人死亡，乙只构成抢劫罪，对丙的死亡结果不承担刑事责任。因为二者没有因果关系，一个人无需对与自己行为没有因果关系的结果负责。（2012年第10题）

例2，甲抢劫丙，将丙打倒在地，在甲知情的情形下，乙参与进来，为了压制反抗，

% Source PDF physical page 157; see audit/source_coverage.json.
踢丙心脏一脚，然后捡起财物。事后查明，丙是被乙踢死的。甲乙构成抢劫罪承继的共同犯罪，乙构成抢劫罪致人死亡，甲也构成抢劫罪致人死亡，因为乙参与进来后，甲和乙已构成共同正犯。

[巩固练习]题1，甲抢劫丙，向丙心脏踢一脚，乙参与进来，为了压制反抗，也向丙心脏踢一脚，后丙死亡。事后无法查明致死一脚是谁踢的。死亡结果只由甲负责。这是因为，如果查明是甲导致的，甲负责，乙不负责。如果查明是乙导致的，乙负责，甲也要负责。如此看来，无论如何，甲都要负责，那么就只让甲负责，乙不需负责。这样不会冤枉甲，如果让乙负责，则有可能冤枉乙。也即，死亡时间无法查明时，死得越早，对大家越有利。

题2，甲入户抢劫，破门入户后，在户内电话通知乙为自己的抢劫望风，乙赶到后在户外望风。甲乙构成抢劫罪的共同犯罪，甲是实行犯，并适用“入户抢劫”的加重法定刑，乙是承继的帮助犯，但不适用“入户抢劫”的加重法定刑，因为当甲入户时，乙没有参与，乙对甲的入户没有帮助，乙只帮助甲在户内抢劫。

\subsection*{二、中途退出：共犯关系的脱离}

有人想中途退出，想成立犯罪中止，条件：脱离共犯关系，具体而言，消除自己的贡献，包括物理性、心理性的贡献。

\InsertDiagram{figures/manual_table_p157_1.tex}

例1（教唆犯），甲教唆乙入室盗窃，乙在去的路上，甲又后悔，便打电话让乙回来，但乙不答应。甲便不管。乙盗窃既遂。教唆犯的贡献是引起他人犯意，因此要成立中止，要求打消对方的犯意。甲没有打消乙的犯意，不能成立犯罪中止，要对乙的既遂结果负责，也构成犯罪既遂。

例2（帮助犯），甲为乙的盗窃提供了一把钥匙，又后悔，向乙索要，乙暗自配了一把（或在甲知情情况下配了一把），然后将原来的钥匙还给甲，用配的钥匙盗窃成功。甲没有切实消除帮助作用，不成立中止，而成立既遂。如果甲要回钥匙，乙用其他方法盗窃既遂，则甲成立中止。（2008年第19题，2011年第55题）

例3（帮助犯），甲答应为乙的盗窃望风，但是又感到害怕，便悄悄溜走。乙在里面不知道甲已经溜走。乙盗窃既遂。甲没有明确告知乙自已要走，没有消除对乙的心理性贡献，不成立中止。乙既遂，甲也既遂。

例4（共同正犯），甲乙共谋次日共同入室盗窃，第二天甲欲中止，对乙谎称：“我胃疼，去不了，你去吧！”乙相信便自己去实施，并盗窃既遂。（1）甲没有明确告知退出意思，没有消除心理性贡献，不成立中止，而成立既遂。（2）如果甲明确告知“自已不想干了”，则可以成立犯罪中止，因为此时尚处在预备阶段，甲的贡献只有心理性贡献。消

% Source PDF physical page 158; see audit/source_coverage.json.
除了心理性贡献，便可以成立犯罪中止。（3）在此，成立犯罪中止，不要求有效防止既遂结果发生。（2022年试题）

[巩固练习］题1（共同正犯，2015年主观题），甲乙共谋杀丙，由甲打电话将丙诱骗到湖边小屋，甲照办。甲又告诉乙都别干了，乙不答应。甲便给丙打电话，想告诉丙赶紧撤离，但甲未打通，便作罢。乙到小屋杀了丙。甲没有消除物理性贡献，不成立中止。乙既遂，甲也既遂。

题2（共同正犯，2022年试题），甲乙共谋杀丙，一起向丙开枪，第一枪都没打中。甲看到丙可怜，不想杀了，劝乙也不要杀了，但乙未答应，甲便离去。在实行阶段，共同正犯已经实施了实行行为。成立中止，要求消除物理性、心理性贡献，是指要求有效防止既遂结果发生。

（1）乙继续开枪，杀死丙，乙既遂。甲要对乙制造的结果负责，甲也既遂。

（2）乙继续开枪，仍未打中，此时被警察抓捕。乙构成未遂。由于没有发生既遂结果，所以甲的中止行为可以成立犯罪中止。\footnote{成立犯罪中止，要求有中止行为，要求既遂结果未发生，但不要求既遂结果未发生是中止行为的功劳。参见“犯罪形态”一讲的“犯罪中止”一节。}

（3）乙继续开枪时，丙极力求饶，乙可怜丙，放弃杀丙。乙构成中止。由于没有发生既遂结果，所以甲的中止行为可以成立犯罪中止。

\OriginalPanel{典型真题}

\ensuremath{>}\ensuremath{>}

甲乙约定好去丙家抢劫。下列情形中，甲构成抢劫罪既遂的有？（2024年回忆版金题）\footnote{[答案]A项，甲没有中止行为，对乙的既遂负责。B项，在预备阶段，甲明确告知乙自己放弃，成立中止。C项，在实行阶段，要成立中止，原则上要求有效防止既遂结果发生。甲不成立中止，对乙的既遂负责。D项，甲实施了阻止行为，虽然原则上要求有效防止既遂结果发生，但是甲被打晕，甲已经尽力了，能够成立中止。本题答案为AC。}

A.甲去丙家的途中，被父母强行拉回家。乙抢劫既遂

B.甲去丙家的途中，心生悔意，自行回家，并告知乙。乙抢劫既遂

C.甲乙共同压制丙的反抗后，甲心生悔意，自行离去。乙继续抢劫，并既遂

D.甲乙共同压制丙的反抗后，甲心生悔意，阻止乙，被乙打晕。乙继续抢劫，并既遂

\section{片面的共同犯罪}

按照正常情况，二人构成共同犯罪，应是相互构成。片面的共同犯罪，是指单方面构成共同犯罪。

\subsection*{一、片面帮助}

这是指甲暗中帮助乙实行犯罪，而乙对此并不知情。例如，甲欲杀丙，看到乙在追杀丙，便暗中设置绳索将丙绊倒（轻微伤），乙顺利杀了丙，乙不知道甲帮了自己。

肯定说认为：二人要成立共同犯罪，不要求相互有意思联络，不要求均有参与意识，只要一方有参与意识即可。有参与意识的一方可以成立单方面的共同犯罪，对另一方的违

% Source PDF physical page 159; see audit/source_coverage.json.
法事实要负责。片面的共同犯罪，就是指单方面的共同犯罪。无参与意识的一方构成独自的单独犯罪。基于此，从甲的角度看，甲构成乙的故意杀人罪的片面的帮助犯，应对乙的杀人的违法事实负责。\CJKunderwave{肯定说的意思是，“我只要暗恋你，就要对你负责，我的行为要与}\CJKunderwave{你的行为挂钩}”。从乙的角度看，乙构成单独的故意杀人罪。

否定说认为：二人要成立共同犯罪，要求二人有相互的参与意识（意思联络）。因此，甲与乙不构成共同犯罪，甲对乙制造的违法事实不用负责，甲只对自己的行为负责。换言之，甲要对乙的违法事实负责，需要乙知道甲在帮乙。\CJKunderwave{否定说的意思是，“因为我不知道}\CJKunderwave{你喜欢我，所以你不需要对我的事负责，你的行为不能与我的行为挂钩”}。所以，甲不构成乙的故意杀人罪的片面的帮助犯，对甲应单独处理。由于甲只是导致丙轻微伤，因此不构成犯罪。

[提示】否定说也感觉自已定无罪的结论不合理，因此当今的否定说也转而承认肯定说的结论。但是，在下文片面实行的场合，否定说仍坚持自己的观点。不过，考试如果考片面帮助犯有哪些观点展示，则仍需要将否定说展示出来。

[注意]片面望风问题。例如，乙看到甲入户盗窃，主动为甲望风，甲对此不知情。期间无异常情况发生。甲盗窃既遂。第一，由于甲不知道乙在望风，所以乙的望风对甲没有心理性帮助作用。第二，乙的望风有无物理性帮助作用？少数说认为，虽然望风期间无异常情况，但是主人有可能回家，因此，望风还是起到了可能的物理性帮助作用。所以，乙构成片面的帮助犯。多数说认为，片面的帮助犯的成立条件比普通的帮助犯应高一些，物理性帮助行为只有发挥实际的物理性帮助作用，才能成立片面的帮助犯。乙的望风没有发挥实际的物理性帮助作用，也无心理性帮助作用，因此不构成片面的帮助犯。（2021年试题，2023年试题）

\subsection*{二、片面教唆}

这是指甲暗中教唆乙犯罪，而乙没有认识到被教唆。例如，甲偷偷将乙的妻子与丙通奸的照片放在乙路过的地方，乙偶然发现，以为是丙不慎弄丢的，火冒三丈，杀了丙。

肯定说认为：甲构成乙的故意杀人罪的片面的教唆犯，应对乙杀人的违法事实负责。从乙的角度看，乙构成单独的故意杀人罪。也即，这里的共同犯罪不是相互构成，而是单向性构成。

否定说认为：甲不构成乙的故意杀人罪的片面的教唆犯，对甲应单独处理，不构成犯罪。由于这种结论不合理，因此当今的否定说也转而承认肯定说的结论。

\subsection*{三、片面实行}

这是指甲暗中和乙共同实行犯罪，而乙对此不知情。在此存在观点展示。肯定说是多数观点。

例如，甲得知乙欲强奸妇女，便提前将妇女打晕（轻伤），然后退出，乙顺利强奸了妇女，但乙不知是甲将妇女打晕。强奸罪的实行行为包括暴力行为和奸淫行为，所以甲属于片面实行。

肯定说（多数说）：（1）甲有帮助乙犯罪的参与意识，应对乙的违法事实（强奸）负责，结合自己的暴力行为，构成乙的强奸罪的片面的实行犯（正犯）。（2）乙构成单独的强奸罪。

% Source PDF physical page 160; see audit/source_coverage.json.
否定说（少数说）：由于甲实施了实行行为，一般而言，能够单独构成犯罪，不会导致无罪的不当结论，因此坚持认为，两人要成立共同犯罪，需要有相互的意思联络。（1)在强奸妇女这件事上，由于乙不知道甲在帮自己，所以甲乙不构成共同实行犯，也即甲不构成乙的强奸罪的片面的实行犯，而构成单独的故意伤害罪。\CJKunderwave{注意}：否定说虽然不承认“片面的实行犯”概念，但依然承认“片面的帮助犯”概念，也即承认甲构成乙的强奸罪的片面的帮助犯。最终结论是，甲同时触犯故意伤害罪和强奸罪的片面帮助犯，想象竞合，择一重罪论处。（2）乙构成单独的强奸罪。

[巩固练习]题1（2017年第54题），甲知道乙计划前往丙家抢劫，为帮助乙取得财物，便暗中先赶到丙家，将丙打昏（轻伤）后离去。乙来到丙家时，发现丙已昏迷，以为是丙疾病发作晕倒，盗窃既遂。

肯定说：（1）分析甲。甲有帮助乙犯罪的参与意识，应对乙的违法事实（盗窃）负责，甲的暴力行为能够与乙的盗窃行为挂钩，基于此，甲便构成抢劫罪，属于片面的实行犯，意思是单方面构成抢劫罪的实行犯。（2）从乙的角度出发，乙是单独犯罪，构成盗窃罪既遂。

否定说：（1）分析甲。在抢劫丙这件事上，由于乙不知道甲在帮自己，所以甲乙不构成共同实行犯，也即甲对乙的盗窃事实不用负责，甲的行为不能与乙的行为挂钩。甲构成单独的故意伤害罪。不过，否定说承认“片面的帮助犯”概念，承认甲构成乙的盗窃罪的片面帮助犯。最终结论是，甲同时触犯故意伤害罪和盗窃罪的片面帮助犯，想象竞合，择一重罪论处。（2）从乙的角度考察，乙是单独犯罪，构成盗窃罪既遂。

题2（2019年试题），乙请甲为自已的盗窃望风，仅要求甲看到主人丙回家就电话告知他即可。乙在户内盗窃时，甲看到丙回家，为了帮助乙顺利取得财物，使用暴力阻拦，将丙打成重伤。乙盗窃既遂，但对外面的事情不知情。

本案与上文题1的不同之处在于，甲乙事先存在共同犯罪（盗窃罪）的基础，甲一开始就构成乙的盗窃罪的帮助犯。该帮助犯不是片面的帮助犯，因为乙知道甲在帮自已望风。由于甲乙有共同犯罪的基础，所以甲多干了一件事（对人暴力），也被称为实行过限。对于甲多干的事，由于乙不知情，故乙不用负责。

肯定说：（1）分析甲。虽然乙不知道甲使用暴力，但是甲实施暴力是为了帮乙，所以，甲的暴力行为能够与乙的盗窃事实挂钩，基于此，甲便构成抢劫罪，属于抢劫罪（致人重伤），属于片面的实行犯。（2）从乙的角度出发，乙不是单独犯罪，而是与甲构成盗窃罪的共同犯罪，乙是实行犯。由于乙对甲多干的事（暴力行为）不知情，故不用负责。最终结论，甲乙在盗窃罪的范围成立共同犯罪，甲最终定抢劫罪（片面实行犯），乙定盗窃罪（实行犯）。

否定说：（1）分析甲。虽然乙知道甲的望风行为，但是不知道甲的暴力行为，所以就甲的暴力行为而言，二人不构成共同犯罪，也即甲的暴力行为不能与乙的盗窃事实挂钩，只有甲的望风行为与乙的盗窃事实挂钩。因此，甲不构成抢劫罪。甲构成单独的故意伤害罪。同时，甲构成乙的盗窃罪的帮助犯。甲同时触犯故意伤害罪和盗窃罪的帮助犯，想象竞合，择一重罪论处。（2）从乙的角度考察，乙与甲构成盗窃罪的共同犯罪，乙是实行犯。乙对甲的暴力行为不用负责。最终结论，甲乙在盗窃罪的范围成立共同犯罪，甲同时触犯故意伤害罪和盗窃罪的帮助犯，想象竞合，择一重罪论处；乙定盗窃罪（实行犯）。

% Source PDF physical page 161; see audit/source_coverage.json.
\section{共同犯罪的结合问题}

\subsection*{一、共同犯罪与身份犯}

\subsubsection*{（一）真正身份犯}

1.角色担当

（1）真正身份犯中有定罪身份，有身份者才有资格构成正犯（包括直接正犯和间接正犯）。例如，只有国家工作人员才可能构成贪污罪的直接正犯、间接正犯。

（2）无身份者不能构成正犯。例如，官员贪污，作为普通公民的妻子不能构成贪污罪的直接正犯和间接正犯。（2013年第55题）

（3）无身份者可以构成共犯（教唆犯、帮助犯）。例如，官员贪污，作为普通公民的妻子可以构成贪污罪的教唆犯、帮助犯。\footnote{无身份者能否构成共同正犯？多数说认为可以。例如，丈夫是官员，妻子策划和主导丈夫的贪污行为，丈夫只是配合演出。丈夫构成贪污罪的正犯。妻子的作用大于普通的教唆犯、帮助犯，但妻子无法构成贪污罪的间接正犯，不过可以构成贪污罪的共同正犯。按照多数说，真正身份犯中，所谓“有身份者才有资格构成正犯”，其中的“正犯”是指单独正犯，不包括共同正犯。}

2.分析路径：分别分析每个人，若同时触犯两个罪，想象竞合，择一重罪论处。

\InsertDiagram{figures/manual_table_p161_1.tex}

例1，妻子帮助丈夫（国家工作人员）收受贿赂。从丈夫角度看，丈夫是受贿罪的实行犯，妻子是共犯（帮助犯）。妻子不构成其他罪的实行犯。

例2，甲乙开车来到某公司，欲盗窃公司仓库，对公司保安经理王某（正式职员）说：“只要你放行，弄到值钱的，有你一份。”王某答应。从王某角度看，王某没有实施职务侵占罪的实行行为，故不构成职务侵占罪。从甲乙角度看，甲乙构成盗窃罪的实行犯，王某构成共犯(帮助犯)。

[提示］早期理论认为，谁有身份就按谁的身份来定罪名。如果这样，那么例2就需要按照王某的身份（公司职员）定职务侵占罪。问题是王某只是帮助者，不可能是职务侵占罪的实行犯，而甲乙又不是公司职员，也不可能构成职务侵占罪的实行犯，而一项共同犯罪不可能缺少实行犯。

例3，甲（投保人）和乙（保险公司职员）合谋诈骗保险公司的保险金，由甲负责制造保险事故、编造虚假索赔材料，由乙负责通过欺骗其他同事，办理保险公司理赔。二人共诈骗保险金100万元。（1）从甲的角度看，甲乙构成保险诈骗罪的共同犯罪，甲是实行

% Source PDF physical page 162; see audit/source_coverage.json.
犯，乙是帮助犯。（2）从乙的角度看，甲乙构成职务侵占罪的共同犯罪，乙是实行犯，甲是帮助犯。（3）甲同时触犯了保险诈骗罪（实行犯）和职务侵占罪（帮助犯），择一重罪论处，定保险诈骗罪的实行犯。乙同时触犯了保险诈骗罪（帮助犯）和职务侵占罪（实行犯），择一重罪论处，定职务侵占罪的实行犯。（4）传统理论认为，甲乙是共同犯罪，应定同一个罪名，例如按照乙的身份来定，都定职务侵占罪。问题是，如此，甲定职务侵占罪的帮助犯，而这比给甲定保险诈骗罪的实行犯处罚要轻，就不当地轻纵了甲。

[注意］司法解释的特殊规定：公司、企业或其他单位中，非国家工作人员与国家工作人员勾结，分别利用各自的职务便利，共同将本单位的财物非法占为己有，按照主犯的犯罪性质定罪。二人只能定同一个罪名。\footnote{2000年7月8日最高人民法院《关于审理贪污、职务侵占案件如何认定共同犯罪几个问题的解释》。}

例如，国有控股公司里，甲是会计（非国家工作人员），乙是财务总监（国家工作人员），共谋利用各自职务便利，共同侵吞公司财物。按照正常原理，从乙的角度看，乙构成贪污罪的实行犯，甲是共犯。从甲的角度看，甲构成职务侵占罪的实行犯，乙是共犯。乙同时触犯贪污罪的实行犯和职务侵占罪的共犯，想象竞合，择一重罪论处，定贪污罪的实行犯。甲同时触犯职务侵占罪的实行犯和贪污罪的共犯，想象竞合，择一重罪论处，职务侵占罪的实行犯可能处罚更重，定职务侵占罪的实行犯。

但是，按照司法解释，二人只能定同一个罪名，例如，都定贪污罪。然而，对甲定贪污罪的共犯，处罚比职务侵占罪的实行犯可能更轻，如此便不当地轻纵了甲。司法解释要求定同一个罪名，背后原因是完全犯罪共同说。实际上，两个人构成共同犯罪，并不要求定同一个罪名。不过，司法解释具有法律效力，考试时按照司法解释来答题。该司法解释只针对“贪污罪+职务侵占罪”的情形。对于其他情形依照正常原理处理。

\subsubsection*{（二）不真正身份犯}

不真正身份犯中有量刑身份。无身份者与有身份者共同犯罪，在量刑时，对无身份者不能适用量刑身份的法定刑。

例如，法条规定，国家机关工作人员犯诬告陷害罪，应从重处罚。国家机关工作人员就是诬告陷害罪的量刑身份。甲（普通公民）和乙（国家机关工作人员）共同诬告陷害王某。甲乙构成诬告陷害罪的共同犯罪，对乙应从重处罚，但对甲不能从重处罚。

\subsection*{二、共同犯罪与不作为犯}

\InsertDiagram{figures/manual_table_p162_1.tex}

% Source PDF physical page 163; see audit/source_coverage.json.
\ContinuedTable{figures/manual_table_p163_1.tex}

\subsection*{三、共同犯罪与实行过限}

实行过限，是指甲乙共同实施A罪，实行犯乙多实施了B行为。对于B行为及结果，甲是否承担刑事责任？判断标准：如果超出共同故意的范围，则不负故意责任，但有可能负过失责任。

实行过限的行为主体主要是实行犯，也称实行过剩，也称共同犯罪的过剩或共犯过剩。实行过限存在于两种类型：（1）共同正犯之间。（2）共犯（教唆犯和帮助犯）+正犯。

\subsubsection*{（一）共同正犯}

例1，甲乙共同盗窃丙家，甲在客厅偷，乙在卧室偷，乙看到睡着的女主人丙，又强奸了丙。对此，甲看到了，没有搭理，继续偷拿财物。甲对乙的强奸罪要不要负刑事责任？乙的强奸罪超出了甲乙共同故意（盗窃罪）的范围，因此甲对此不用负责。\footnote{完整的分析方法是，先分析客观要件，再分析主观要件。不过，从做题的效率角度看，用简单的标准只要能把题做对就行。以后从事实务，分析案件，则需要从客观要件到主观要件，一层一层分析论证。考试做题与实务分析案件还是有区别的。}

例2（2010年第93题），甲乙共同诈骗到丙的财物，丙抓捕，乙为了抗拒抓捕，将丙打成重伤，但甲对此明确反对。乙构成事后转化型抢劫罪。乙的暴力行为超出了甲乙共同故意（诈骗罪）的范围，因此甲对此不用负责。甲只负诈骗罪的责任。

[提示1]在此常见错误看法是，根据是否知情来判断；甲对乙的过限事实如果不知情，则不负责；如果知情，则要负责。正确看法是，不知情，肯定不用负责，但知情，不一定就要负责，还要看有无犯罪故意。

[提示2】只有共同正犯，没有教唆犯和帮助犯，这种共同犯罪被称为简单共同犯罪。“正犯+教唆犯或帮助犯”，这种共同犯罪被称为复杂共同犯罪。这些共同犯罪中，实行犯都可能出现实行过限。例如（2023年试题），“只有复杂共同犯罪（有教唆犯和帮助犯）中，才存在实行过剩”。该说法是否正确？答案：说法错误。例如，上述例1、例2中，甲乙属于简单共同犯罪，但存在实行过剩（实行过限）。

% Source PDF physical page 164; see audit/source_coverage.json.
\subsubsection*{（二）正犯+教唆犯}

例1，甲教唆乙盗窃丙家，乙在盗窃时，看到女主人丙颇有姿色，便强奸了丙。乙的强奸罪超出了甲乙共同故意（盗窃罪）的范围，因此甲对此不用负责。

例2，甲教唆乙盗窃丙家，乙盗窃得手后，丙发现，抓捕乙，乙为了抗拒抓捕，将丙打成重伤。乙构成转化型抢劫罪（既遂）。乙的暴力行为超出了甲乙共同故意（盗窃罪）的范围，因此甲对此不用负责。甲只负盗窃罪（既遂）的责任。

例3（2007年第60题），甲教唆乙抢夺丙，乙抢夺得手后，丙抓捕乙，乙为了抗拒抓捕，将丙打成重伤。乙构成转化型抢劫罪（既遂）。乙的暴力行为超出了甲乙共同故意（抢夺罪）的范围，因此甲对此不用负责。甲只负抢夺罪（既遂）的责任。

[巩固练习]题1，甲雇佣乙伤害丙，告诫不要闹出人命。\footnote{注意，雇佣者是教唆犯，不是间接正犯，因为花钱雇佣人，对人没有刑法上的支配力。}乙在伤害过程中，由于丙激烈反抗，乙便杀死了丙。乙构成故意杀人罪既遂。乙的杀人行为超出了甲乙共同故意（只伤害，不杀人）的范围，因此甲对杀人不负故意犯罪的责任，但应负过失犯罪的责任，因为甲对杀人有过于自信的过失。因此，甲不构成故意杀人罪，而构成故意伤害罪（过失）致人死亡。\footnote{此时，若表述为甲构成故意伤害罪（过失）致人死亡（教唆犯），则违反一个结论：教唆犯的罪名不能是过失犯罪。准确的表述应是，甲构成故意伤害罪（教唆犯），量刑按照故意伤害罪中的“过失致人死亡的，加重处罚”处理。换言之，甲对过失致人死亡负责时，甲的角色不是教唆犯，而是正犯。也即，在致人死亡上，甲是正犯，乙也是正犯，区别在于甲是过失致人死亡罪的正犯，乙是故意杀人罪的正犯。甲对乙制造的死亡结果负责，是因为甲的雇佣行为对乙的杀人起到促进作用。}

题2，甲教唆乙使用暴力抢劫丙，但告诫不要闹出人命。乙使用暴力抢劫丙时，丙奋力反抗，乙杀死丙，抢劫到财物。乙构成抢劫罪（故意）致人死亡。乙的杀人行为超出了甲乙共同故意（只抢财物，不杀人）的范围，因此甲对杀人不负故意犯罪的责任，但应负过失犯罪的责任，因为甲对杀人有过于自信的过失。甲构成抢劫罪（过失）致人死亡。

\subsubsection*{（三）正犯+帮助犯}

例如，甲帮助乙盗窃丙家，在外望风。乙在屋内盗窃时，看到女主人丙颇有姿色，便强奸了丙。乙的强奸罪超出了甲乙共同故意（盗窃罪）的范围，因此甲对此不用负责。

[巩固练习]题1（2011年第55题），甲帮助乙盗窃丙家，在外望风，乙在屋内盗窃时被主人丙发现，为了继续取得财物，将丙打晕，取得财物。乙构成抢劫罪实行犯（既遂）。乙的抢劫罪超出了甲乙共同故意（盗窃罪）的范围，因此甲对此不用负责。客观上，由于乙的抢劫行为可以包容评价为盗窃行为，因此甲有帮助盗窃的行为。主观上，甲有帮助盗窃的故意。因此，甲构成盗窃罪的帮助犯（既遂）。

[提示]如果认为抢劫罪与盗窃罪是公开与秘密的对立关系，则抢劫罪无法包容评价为盗窃罪，由此导致对甲无法定盗窃罪的帮助犯（既遂）。这显然是不妥当的。有同学认为，即使乙的抢劫无法包容评价为盗窃，给甲也能定盗窃罪的帮助犯。但是，这种看法违反了共犯从属性原则。如果实行犯不构成盗窃罪（既遂），则帮助犯不可能构成盗窃罪的帮助犯（既遂）。

题2，乙入户伤害丙，甲在外望风，告诫不要闹出人命。乙在伤害过程中，由于丙激烈反抗，乙便杀死了丙。乙构成故意杀人罪既遂。乙的杀人行为超出了甲乙共同故意（只

% Source PDF physical page 165; see audit/source_coverage.json.
伤害，不杀人）的范围，因此甲对杀人不负故意犯罪的责任，但应负过失犯罪的责任，因为甲对杀人有过于自信的过失。因此，甲不构成故意杀人罪，而构成故意伤害罪致人死亡。（2008年四川第61题引申情形）

\subsubsection*{（四）分析完作为，再分析不作为}

以上分析，只是从作为的角度考察了甲对乙的过限事实要不要负刑事责任。若从作为角度得出肯定结论，则不需要考察不作为角度。若作为角度得出否定结论，则需要考察不作为角度。

例1，甲乙共同盗窃丙家，甲在客厅偷，乙在卧室偷，乙看到睡着的女主人丙，又强奸了丙。对此，甲看到了，没有搭理，继续偷拿财物。甲对乙的强奸罪要不要负责？从作为的角度看，甲不用负责，因为甲没有实施强奸的共同故意。接下来从不作为的角度看，甲对乙的强奸罪没有阻止义务，因为甲的先行行为只是通过平和手段盗窃财物，没有使丙陷入危险状态。因此，甲不构成强奸罪的不作为的帮助犯。最终，甲对乙的强奸罪不负刑事责任。

例2，甲乙共同抢劫丙女，共同将丙打晕，甲在丙身上取财，乙又强奸丙，甲看到了，不阻止，只是骂了句：“能不能专注点？”甲对乙的强奸罪要不要负责？从作为的角度看，甲不用负责，因为甲没有实施强奸的共同故意。接下来从不作为的角度看，甲对乙的强奸罪具有阻止义务，因为甲的先行行为（打晕丙）导致丙陷入危险状态。因此，申构成强奸罪的不作为的帮助犯。许多同学容易混淆这两个角度，需要注意。当然，一般情况下不需要特别考虑不作为的角度，除非题干强烈暗示了这一点。

\subsubsection*{（五）实行不足}

实行不足，是指预定的共同犯罪是重罪，实行犯只实施了轻罪。其他共犯人对轻罪负刑事责任。这里的轻罪与重罪是性质相同、程度不同的两个罪。重罪可以包容评价为轻罪。

例1（2011年第55题），甲在外望风，帮助乙杀人，乙在屋里仅实施了伤害。客观上，甲对乙的伤害行为提供了帮助。主观上，甲有帮助杀人的故意，可以包容评价为帮助伤害的故意。因此，甲构成故意伤害罪的帮助犯。如果甲帮助乙伤害，乙却实施了杀人，则属于实行过限的问题。

例2，甲教唆乙抢劫，乙仅实施了盗窃。甲有教唆抢劫的故意，可以包容评价为甲有教唆盗窃的故意。因此，甲构成盗窃罪的教唆犯既遂。如果乙仅实施了强制猥亵罪，则属于实行过限的问题，甲对强制猥亵罪不用负责。

[巩固练习]题1（2011年第55题），甲教唆乙入户抢劫，乙仅实施了普通抢劫。情节加重犯可以包容评价为基本犯，也即入户抢劫可以包容评价为抢劫。因此，甲教唆乙入户抢劫，可以包容评价为甲教唆乙普通抢劫。甲构成普通抢劫罪的教唆犯。

题2，甲教唆乙实施抢劫罪，乙实施了入户抢劫。甲构成抢劫罪的教唆犯，对于乙的入户，由于没有教唆故意，因此不构成入户抢劫的教唆犯，也不构成非法侵入住宅罪的教唆犯。对甲不能适用“入户抢劫”的升格法定刑。这种情形不属于实行不足，而属于实行过限。

[结论］性质相同、程度不同的重罪与轻罪。（1）甲教唆或帮助乙实施重罪，乙只实

% Source PDF physical page 166; see audit/source_coverage.json.
施了轻罪，甲成立轻罪的教唆犯、帮助犯。（2）甲教唆或帮助乙实施轻罪，乙实施了重罪，甲成立轻罪的教唆犯、帮助犯。

\subsection*{四、共同犯罪与认识错误}

共同犯罪与认识错误的结合问题，需要用到前文关于事实认识错误的知识、实行过限的知识。先辨认认识错误的类别，再进行处理。

\subsubsection*{（一）共同正犯}

例如（2004年第18题），甲、乙共谋杀害在博物馆工作的丙，两人潜入博物馆同时向丙各开一枪，甲击中丙身边的珍贵文物，造成文物毁损；乙未击中任何对象。甲存在不同构成要件间的打击错误。甲一方面构成故意杀人罪未遂，另一方面构成过失损毁文物罪，想象竞合，择一重罪论处，定故意杀人罪未遂。乙构成故意杀人罪未遂。

\subsubsection*{（二）教唆犯或帮助犯+正犯}

例如（实行犯是对象错误+教唆犯是打击错误），甲教唆乙杀害丙，将丙的照片给乙，提醒乙不要认错人、杀错人，乙在丙下班的路上守候，误将丁当作丙而杀害。（1）乙构成对象错误。乙对丙构成故意杀人罪犯罪预备，对丁构成故意杀人罪既遂，想象竞合，择一重罪论处，定故意杀人罪既遂。（2）实行犯乙产生对象错误，教唆犯甲构成打击错误。依照具体符合说，甲对丙构成故意杀人罪犯罪预备，对丁构成过失致人死亡罪，想象竞合，择一重罪论处。依照法定符合说，甲对丙构成故意杀人罪犯罪预备，对丁构成故意杀人罪既遂，想象竞合，择一重罪论处。

[巩固练习]题1（实行犯没有错误+教唆犯是打击错误），甲教唆乙杀丙，言明不要造成其他后果。乙几次杀丙均未成功，后来对丙的住宅（周边没有其他人与物）进行爆炸，结果只是将丙的妻子丁炸死。（1）乙没有认识错误，对丁的死亡至少持放任态度，对丁构成故意杀人罪既遂，对丙构成故意杀人罪未遂，想象竞合，择一重罪论处，定故意杀人罪既遂。（2）甲对丁的死亡至少持过于自信的过失，基于此甲构成打击错误。根据具体符合说，甲对丙构成故意杀人罪未遂，对丁构成过失致人死亡罪，想象竞合，择一重罪论处，定故意杀人罪未遂。根据法定符合说，甲对丙构成故意杀人罪未遂，对丁构成故意杀人罪既遂，想象竞合，择一重罪论处，定故意杀人罪既遂。（2008年第55题）

题2（实行犯没有错误+教唆犯是对象错误），甲欲杀害丙，误以为前方的丁是丙，手指向丁，指示杀手乙：“将那个人开枪打死！”乙照办，丁死亡。由于乙没有辨认对象身份的任务，只有机械执行的任务，故不存在对象认识错误。甲存在对象错误。乙构成故意杀人罪既遂，甲也构成故意杀人罪既遂。（2022年试题）

\subsubsection*{（三）犯罪人成为被害人}

这是指，正犯产生对象错误，导致共犯人（教唆犯、帮助犯）、正犯自己、其他正犯成为被害人。这类案件可称为“乌龙案”。

1.共犯人成为被害人

例1（渔网案，2019年试题），甲欲盗窃丙的渔网，请乙提供工具，乙答应。甲拿着乙的工具盗窃时，误将乙的渔网当作丙的渔网而盗窃了。丙的渔网当时在现场，但是甲没有注意到。

% Source PDF physical page 167; see audit/source_coverage.json.
（1）分析甲。第一，针对预定目标丙的渔网，盗窃渔网的着手，应当是手即将碰到渔网，由于甲没有注意到，没有着手实行，所以甲构成盗窃罪犯罪预备。第二，针对实害对象乙的渔网，甲存在对象错误，但不影响盗窃罪既遂，因为盗窃罪只要求盗窃“他人”的财物就行。最终，甲整体上的一个行为同时触犯盗窃罪犯罪预备和盗窃罪既遂，想象竞合，择一重罪论处，定盗窃罪既遂。

（2）分析乙。第一，针对预定目标丙的渔网，由于实行犯甲构成盗窃罪犯罪预备，因此乙构成盗窃罪帮助犯（犯罪预备）。第二，针对实害对象乙的渔网，乙不用负责，因为乙是被害人。最终，乙仅构成盗窃罪帮助犯（犯罪预备）。\footnote{当实行犯甲存在对象错误，帮助犯乙构成打击错误，因为乙对实害对象及结果的心理是过失或意外事件。对于打击错误，存在具体符合说和法定符合说。但两种学说都认为：针对预定目标丙的渔网，由于实行犯甲构成盗窃罪犯罪预备，因此乙构成盗窃罪帮助犯（犯罪预备）；针对实害对象乙的渔网，乙不用负责。所以，在此不需要进行观点展示。}

例2，乙女教唆甲男强奸室友丙女。某晚，甲男入室，丙女恰巧不在，甲误将乙女当作丙女强奸了。

（1）分析甲。第一，针对预定目标丙，甲已经有了预备行为，未能着手，构成强奸罪犯罪预备。第二，针对实害对象乙，甲存在对象错误，但不影响强奸罪既遂。最终，甲整体上的一个行为同时触犯强奸罪犯罪预备和强奸罪既遂，想象竞合，择一重罪论处，定强奸罪既遂。

（2）分析乙。第一，针对预定目标丙，乙构成强奸罪教唆犯（犯罪预备）。第二，针对实害对象乙，乙不用负责，因为乙是被害人。最终，乙仅构成强奸罪教唆犯（犯罪预备）。提示：许多同学认为乙无罪。这种说法不准确。

2.“教唆犯或帮助犯+正犯”中，正犯自己成为被害人

例如，乙参加甲的婚宴，中途要离开，甲送乙。在门口，乙找不到自己的摩托车。有点醉酒的甲对乙讲：“你给我拿个工具，我给你撬一辆。”乙将工具给甲。甲拍着一排摩托车中的一辆，对乙说：“我给你撬这辆！”说着就撬。乙说：“那辆不错，撬那辆。”甲照办，撬开，让乙骑走。次日，甲发现撬的是自已的摩托车。乙方知撬的是甲的车。

（1）分析实行者甲。第一，甲乙的预定目标是他人的车。甲对他人的车已经着手，构成实行阶段的中止。第二，针对实害对象甲的车，甲是被害人，甲不构成犯罪。\footnote{本案与“行为对象的转移”（“罪数”一讲中“两个行为”一节）有所不同。后者中，对第一个盗窃对象不定罪，是因为对第二个盗窃对象能定罪。而本案中，对第二个盗窃对象（甲的车）不能定罪，所以对第一个盗窃对象需要定罪。}

（2）分析帮助者乙。第一，针对预定目标他人的车，乙构成犯罪中止。第二，针对实害对象甲的车，根据共犯从属性原理，由于实行者甲不构成犯罪，所以帮助者乙也不构成犯罪。

3.共同正犯中，另一正犯成为被害人

例如，甲乙共同入户抢劫丙家，约定抢劫后如果丙追赶，就都向丙开枪。二人抢劫后，甲发现丙追赶，就开枪，实际是乙，打伤乙。乙当时正向外逃跑。丙当时在现场，甲没有看到丙。

（1）分析正犯甲。第一，甲对预定目标丙构成故意杀人罪犯罪预备。第二，甲对实害对象乙有对象错误，构成故意杀人罪未遂。想象竞合，择一重罪论处。

% Source PDF physical page 168; see audit/source_coverage.json.
（2）分析正犯乙。第一，由于甲对丙构成犯罪预备，所以乙对丙也构成犯罪预备。第二，实害对象是乙自己，对此乙无罪。

[提示］如果乙当时正要开枪打丙，甲却打中乙，客观上阻止了乙杀丙，那么该案就属于“猪队友帮倒忙”案，是打击错误与偶然防卫相结合的案件。对此分析，参见前面第九讲中“防卫认识”。

\section{共同犯罪的处罚规定}

分析共同犯罪的角色，一般用分工分类法，也即实行犯、教唆犯、帮助犯。在最终处罚时，一般用作用分类法，也即主犯、从犯、胁从犯。

第一，实行犯根据作用大小，可以成为主犯、从犯或胁从犯。

第二，教唆犯根据作用大小，可以成为主犯、从犯或胁从犯。

例1，黑社会老大教唆小弟盗窃，老大是教唆犯，也是主犯。

例2，甲胁迫乙：“你教唆丙去盗窃，你如果不答应，就将你的丑闻曝光！”乙被迫答应。乙教唆丙盗窃。丙便实施了盗窃。乙是教唆犯，也是胁从犯。甲属于间接教唆者，也属于教唆犯。

第三，帮助犯根据作用大小，可以成为从犯或胁从犯，但不可能成为主犯。

[提示］这表明，作用分类与分工分类是从不同角度对共犯人进行划分，二者不是对立关系，而是可以竞合。

\subsection*{一、主犯}

第26条第1款\CJKunderwave{组织、领导犯罪集团进行犯罪活动的}或者\CJKunderwave{在共同犯罪中起主要作用}\CJKunderwave{的}，是主犯。

1.主犯的分类：（1）组织、领导犯罪集团进行犯罪活动的犯罪分子，也就是犯罪集团的首要分子；（2）在其他共同犯罪中起主要作用的犯罪分子。

2.主犯与首要分子的关系

第97条本法所称首要分子，是指在\CJKunderwave{犯罪集团}或者\CJKunderwave{聚众犯罪}中起组织、策划、指挥作用的犯罪分子。

（1）首要分子分为：一是犯罪集团中的首要分子；二是聚众犯罪中的首要分子。

（2）犯罪集团的主犯包括首要分子和其他起主要作用的犯罪分子。因此，犯罪集团的首要分子一定是主犯，但犯罪集团的主犯不一定是首要分子。

（3）聚众犯罪有两种情形：一是只处罚首要分子的聚众犯罪，例如，刑法第291条规定的聚众扰乱公共场所秩序、交通秩序罪。此时如果首要分子只有一人，则无共同犯罪而言（制造违法事实阶段的共同犯罪事实是存在的）。所以，\CJKunderwave{聚众犯罪不一定都是共同犯罪}。二是既处罚首要分子，也处罚积极参与者的聚众犯罪，例如，刑法第290条第1款规定的聚众扰乱社会秩序罪。此时，聚众犯罪属于共同犯罪。

当聚众犯罪中的首要分子有多人时，那么起主要作用的是主犯，起次要作用的是从

% Source PDF physical page 169; see audit/source_coverage.json.
犯。因此，聚众犯罪中的首要分子不一定都是主犯。反过来，聚众犯罪中的主犯也不一定都是首要分子。因为聚众犯罪中其他参与者如果起主要作用，也可以是主犯。

[总结]

\InsertDiagram{figures/manual_p169_1.tex}

[结论】主犯与首要分子的关系：除了\CJKunderwave{犯罪集团的首要分子一定是主犯}外，其他情形下的首要分子和主犯都没有必然的一一对应关系。

3.主犯的处罚

（1）犯罪集团中的首要分子

第26条第3款对组织、领导犯罪集团的首要分子，按照\CJKunderwave{集团所犯的全部罪行}处罚。

[注意】是按照“集团”所犯的全部罪行，而不是按照“全体成员”所犯的全部罪行处罚。

（2）其他主犯

第26条第4款对于第三款规定以外的主犯，\CJKunderwave{应当按照其所参与的或者组织、指挥}\CJKunderwave{的全部犯罪处罚}。（2019年试题）

\subsection*{二、从犯}

第27条在共同犯罪中起\CJKunderwave{次要}或者\CJKunderwave{辅助}作用的，是从犯。

对于从犯，应当从轻、减轻处罚或者免除处罚。

\CJKunderwave{主犯与从犯的关系}：

1.共同犯罪中，可以只有主犯（须两个主犯），没有从犯；但不可能只有从犯，没有主犯。

2.刑法没有规定对从犯的处罚必须“比照”主犯，所以对从犯的处罚可以（但不是应当）轻于主犯。

3.对从犯也需要按照其所参与的全部犯罪处罚。这一点和主犯相同。

\subsection*{三、胁从犯}

第28条对于被胁迫参加犯罪的，应当按照他的犯罪情节减轻处罚或者免除处罚。

胁从犯，是指被胁迫参加犯罪的人。

1.行为人一开始被胁迫参加犯罪，但在着手实行后变得积极主动，在共同犯罪中起主要作用的，应认定为主犯，而非胁从犯。

2.如果行为人身体完全被强制或意志自由完全被剥夺，则不构成胁从犯。例如，劫机分子甲持枪胁迫机长将飞机开往指定地点。机长被迫照办。机长的意志自由完全被剥夺，不构成胁从犯，符合紧急避险条件的，可以成立紧急避险。（2022年试题）

% Source PDF physical page 170; see audit/source_coverage.json.
\InsertDiagram{figures/manual_p170_1.tex}
